\documentclass[11pt,a4paper]{article}
\usepackage[a4paper,margin=23mm,headheight=14pt]{geometry}
\usepackage{kotex}
\usepackage{amsmath,amssymb}
\usepackage{longtable,array}
\usepackage{xcolor}
\usepackage[colorlinks=true,linkcolor=blue!55!black,urlcolor=blue!55!black]{hyperref}
\setlength{\parindent}{0pt}
\setlength{\parskip}{0.55em}
\setcounter{tocdepth}{2}
\renewcommand{\arraystretch}{1.25}
\begin{document}
\begin{titlepage}
\centering
\vspace*{3cm}
{\Large 주식·채권·파생금융상품 3: 실증\par}
\vspace{1.2cm}
{\Huge\bfseries 슬라이드별 학습 가이드\par}
\vspace{1.2cm}
{\large 계약의 현금흐름에서 무차익 가격과 위험관리까지\par}
\vfill
{\large 2026년 2학기\par}
\end{titlepage}
\tableofcontents
\clearpage

\section{01. 파생상품을 공부하는 언어}
이 장의 출발점은 상품 이름이 아니라 계약의 현금흐름이다. 기초자산 가격을 $S_t$, 만기를 $T$로 두고, 누가 언제 무엇을 받는지 식과 그림으로 표현하면 선도·선물·옵션·스왑을 같은 틀에서 이해할 수 있다.\par

\subsection{슬라이드 01 · 한 계약을 보는 세 가지 질문}
파생상품은 어떤 기초자산이나 관측 변수의 값에 따라 미래 지급액이 달라지는 계약이다. 처음부터 세 가지를 적자. \textbf{기초변수는 무엇인가? 지급 시점은 언제인가? 지급액은 그 변수의 어떤 함수인가?} 예컨대 원/달러 선도는 미래 환율 $S_T$에 따라 원화 손익이 변하고, 유럽형 콜옵션은 $\max(S_T-K,0)$처럼 손익이 꺾인다. 이름만 기억하면 계약 규모, 방향, 만기를 바꾼 순간 판단이 어려워진다.\par

기초자산의 실물 인도가 있는 계약도 있고, 차액만 현금으로 결제하는 계약도 있다. 두 방식은 계약의 경제적 노출이 같을 수 있지만 결제 과정에서 필요한 현금과 신용위험은 다르다.\par

\subsection{슬라이드 02 · 시장을 이해하는 학제적 시각}
파생상품은 자산가격이론, 금리의 기간구조, 위험관리와 거래제도가 만나는 영역이다. 미시경제학의 \textbf{한계적 교환}은 위험을 누구에게 이전할지 이해하는 데, 거시경제학의 금리·환율은 계약의 기초변수와 할인율을 이해하는 데 쓰인다. 특히 가격 $F$를 구할 때는 미래 환율이나 주가를 예언하기보다, \textbf{동일한 미래 현금흐름을 만드는 두 거래의 오늘 비용이 같아야 한다}는 무차익 논리를 먼저 사용한다.\par

따라서 계산을 시작할 때는 경제적 질문을 한 줄로 적는 것이 좋다. “미래 1달러를 확정적으로 확보하려면 오늘 얼마가 드는가?” 같은 질문이 선도가격 식을 이끈다.\par

\subsection{슬라이드 03 · 읽을 거리와 학습 순서}
먼저 선도·선물의 결제 구조와 포지션 손익을 익힌 뒤 무차익 가격을 유도하고, 그 다음 옵션·스왑의 비선형성 및 위험관리를 연결하는 순서가 효율적이다. 여러 교재에서 사용하는 $S_0$는 오늘 현물가격, $F_{0,T}$는 오늘 체결하는 만기 $T$의 선도가격, $f_t$는 기존 계약의 시점 $t$ 현재가치다. \textbf{가격 $F$와 가치 $f$는 서로 다른 대상}이다. 새 선도계약은 보통 $f_0=0$이지만 $F_{0,T}$가 0이라는 뜻은 아니다.\par

\par\medskip\noindent\textbf{심화: 교재를 읽을 때 생기는 표기 차이}\quad 어떤 책은 선도 인도가격을 $K$라고 쓰고, 어떤 책은 $F_0$라고 쓴다. 일단 계약을 체결하면 약정 인도가격 $K$는 고정된다. 반면 같은 만기의 새 계약에 적용되는 선도가격 $F_t$는 시점마다 변한다. 이 구별이 기존 선도의 평가식 $f_t=e^{-r(T-t)}(F_t-K)$를 이해하는 열쇠다.\par

\subsection{슬라이드 04 · 평가식보다 중요한 것은 현금흐름 관리}
파생상품을 공부할 때 평가점수나 암기량보다 재현 가능한 풀이 순서가 중요하다. 각 문제마다 (1) 포지션의 방향을 적고, (2) 시점별 현금흐름 표를 만들고, (3) 수량과 통화를 맞추고, (4) 동일 현금흐름의 가격을 비교하자. 이 네 단계로 대부분의 부호 실수를 잡을 수 있다. 특히 $1$계약인지, $1$주인지, 지수 $1$포인트인지 단위를 먼저 명시해야 한다.\par

\subsection{슬라이드 05 · 네 가지 기본 블록과 위험·수익의 교환}
\textbf{선도}는 오늘 정한 가격으로 미래에 사고파는 양자 계약이다. \textbf{선물}은 표준화된 선도를 거래소·청산소 체계에서 매일 정산하는 계약이다. \textbf{옵션}은 정해진 거래를 할 \textbf{권리}를 주되 행사 의무는 주지 않아 비선형 지급액을 만든다. \textbf{스왑}은 여러 날짜의 현금흐름을 교환한다. 예컨대 고정금리와 변동금리를 바꾸는 금리스왑은 여러 기간의 선도금리계약을 묶어 생각할 수 있다.\par

같은 계약은 한 기업에는 보험, 다른 투자자에게는 투기일 수 있다. 환율 상승이 걱정되는 달러 수입업자는 달러 매수 선도로 원화 비용을 고정한다. 반대편의 달러 수출업자는 달러 매도 선도로 원화 수입을 고정한다. 계약의 사회적 역할은 상품명 자체보다 \textbf{기존 노출과 결합한 총 현금흐름}에서 결정된다.\par

\subsection{슬라이드 06 · 지수·로그·복리의 최소 도구}
복리의 핵심은 시간을 더하면 성장배수가 곱해진다는 것이다. 연속복리율 $r$에서 현재 $W_0$는 $T$년 뒤 $W_T=W_0e^{rT}$가 된다. 역으로 미래의 확정금액 $X$의 현재가치는 $Xe^{-rT}$이다. 이유는 오늘 $Xe^{-rT}$를 같은 금리로 운용하면 만기에 정확히 $X$가 되기 때문이다.\par

\[
e^a e^b=e^{a+b},\qquad \frac{e^a}{e^b}=e^{a-b},\qquad \ln(ab)=\ln a+\ln b.
\]

$\ln(W_T/W_0)=rT$이므로 $r=\ln(W_T/W_0)/T$다. 이 식은 뒤에서 환율의 로그변화가 정규분포를 따른다고 가정할 때 다시 쓰인다. 로그는 비율을 합으로 바꾸므로 여러 기간의 수익률을 다루기 편하다.\par

\par\medskip\noindent\textbf{심화: 선도금리 $f_{n,m}$의 뜻}\quad $n$년부터 $m$년까지 적용될 금리를 오늘 정한 값이다. 연속복리 현물금리 $R_n,R_m$이 주어지면 $e^{R_m m}=e^{R_n n}e^{f_{n,m}(m-n)}$이므로 $f_{n,m}=(R_m m-R_n n)/(m-n)$이다. $n$년짜리 투자에 이어 미래 구간을 운용하는 경로와 처음부터 $m$년 운용하는 경로의 만기 금액을 같게 놓은 결과다.\par

\subsection{슬라이드 07 · 다음 장으로 이어지는 분류 기준}
파생상품을 분류할 때는 거래소/장외, 선형/비선형, 기초자산, 결제 방식, 만기, 거래상대방 위험을 따로 살펴야 한다. “주가를 기초로 한다”는 정보만으로 손익의 모양은 정해지지 않는다. 주식 선도는 $S_T-K$로 직선형이지만 주식 콜옵션은 $\max(S_T-K,0)$로 꺾인 선이다. 다음 장에서는 바로 이 \textbf{지급액 함수}부터 시장의 이용 목적까지 연결한다.\par

\section{02. 파생상품의 구조와 시장 참여자}
한쪽의 이익이 다른 쪽의 손실이라는 말은 계약의 지급액을 설명하지만, 거래의 목적까지 설명하지는 않는다. 기존 사업 위험을 줄이는 사람과 그 위험을 받아 수익을 노리는 사람이 같은 계약의 양쪽에 설 수 있다.\par

\subsection{슬라이드 01 · 기초 개념의 지도}
모든 계약은 기초변수 $X$, 만기 $T$, 지급함수 $g$의 조합으로 볼 수 있다. 만기 지급액이 $g(X_T)$라면 선도는 $g(S_T)=S_T-K$, 콜옵션은 $g(S_T)=\max(S_T-K,0)$이다. 이 표기는 가격을 계산하기 전에 계약이 어느 방향으로 위험에 노출되는지 보여 준다.\par

여기서 \textbf{지급액}은 만기에 계약이 주는 돈이고 \textbf{순이익}은 계약을 얻는 데 든 오늘의 비용까지 반영한 결과다. 공정한 새 선도는 보통 처음 지급액이 0이므로 만기의 $S_T-K$가 계약의 만기 손익과 같다. 콜옵션 매수자는 처음에 프리미엄 $C_0$를 지급한다. 금리를 무시하는 단순 예에서 만기 순이익은 $\max(S_T-K,0)-C_0$다. $S_T>K$라서 옵션을 행사해도 $S_T-K<C_0$이면 투자 전체로는 손실이다. 이 차이를 먼저 알아야 나중에 ‘옵션을 행사했다’와 ‘돈을 벌었다’를 혼동하지 않는다.\par

\subsection{슬라이드 02 · 정의에서 용도까지의 연결}
먼저 계약의 지급액을 파악하고, 거래·결제 구조를 살핀 뒤, 헤지·투기·차익거래 중 어떤 목적으로 사용하는지 판단한다. 같은 원/달러 선도 매수라도 달러 지급 의무를 이미 진 수입업자에게는 헤지, 아무 달러 부채가 없는 투자자에게는 환율 상승에 대한 투기가 된다. 거래 목적은 계약 하나가 아니라 \textbf{기존 포지션과 합친 손익}으로 판정한다.\par

예를 들어 3개월 뒤 100만 달러를 지급해야 하는 수입업자에게 환율 $S_T$가 1원 오르면 원화 지급액은 100만 원 증가한다. 100만 달러 매수 선도는 같은 환율 상승에 100만 원 이익을 주므로 기존 노출을 상쇄한다. 반면 달러 지급 의무가 없는 투자자가 똑같은 선도를 사면 환율 상승 이익과 하락 손실을 새로 만드는 셈이다. \textbf{계약의 법적 형태가 같아도 포트폴리오에 더했을 때의 기울기가 다르다.}\par

\subsection{슬라이드 03 · ‘파생’과 제로 순공급의 정확한 뜻}
파생상품의 지급액은 주식, 금리, 환율, 원자재 가격 등 더 기초적인 변수에서 정해진다. 같은 선도계약의 매수자와 매도자를 함께 보면 만기 지급액은 $(S_T-K)+(K-S_T)=0$이다. 이 의미에서 계약의 순공급과 계약 당사자 간 지급액 합계는 0이다. 다만 거래 수수료, 거래상대방 부도, 담보 비용, 세금, 기초자산을 이용한 생산 활동까지 포함한 \textbf{경제 전체의 후생}이 반드시 0이라는 뜻은 아니다. 서로 다른 위험을 가진 두 주체가 위험을 교환하면 양쪽 모두에게 가치가 있을 수 있다.\par

위험을 싫어하는 사람이 미래 소득의 변동을 줄이면 평균 소득이 그대로여도 만족도가 높아질 수 있다. 미시경제학의 오목한 효용함수 $u$에서는 일반적으로 $E[u(W)]\le u(E[W])$다. 수출기업의 달러 매출과 수입기업의 달러 구매비용은 환율에 반대 방향으로 노출된다. 서로 선도를 맺어 각각의 현금흐름 변동을 줄이면, \textbf{계약 지급액의 합계는 여전히 0}이어도 두 기업이 원하는 위험구조에 더 가까워질 수 있다. 물론 거래가격·신용위험·수수료가 너무 불리하면 어느 한쪽에 이득이 없을 수도 있다.\par

\subsection{슬라이드 04 · 지급액 함수를 수학적으로 읽기}
파생상품의 만기 지급액은 현재까지의 기초변수 경로 $\{S_u:0\leq u\leq T\}$의 함수로 쓸 수 있다. 선도는 최종 가격만 쓰는 $S_T-K$이므로 \textbf{경로 독립적}이다. 반면 장벽옵션은 중간에 장벽 $H$를 한 번이라도 넘었는지를 포함하므로 같은 $S_T$에서도 지급액이 달라질 수 있다. 선형 계약은 $S_T$가 1만큼 바뀔 때 손익이 일정하게 바뀌지만, 옵션처럼 꺾이는 계약은 구간마다 민감도가 다르다.\par

\[
\Pi_T^{\mathrm{forward}}=q(S_T-K),\qquad
\Pi_T^{\mathrm{call}}=q\max(S_T-K,0).
\]

이 식의 $q>0$은 매수 수량이다. 선도는 $S_T$가 $K$보다 낮아져도 계약을 이행해야 하므로 손익의 기울기가 계속 $q$다. 콜옵션 매수자는 $S_T\le K$일 때 행사하지 않아 지급액이 0이고, $S_T>K$일 때 기울기가 $q$다. 따라서 선도에는 \textbf{의무}가 있고 옵션에는 \textbf{선택권}이 있다. 이 선택권이 일방적으로 유리하므로 콜을 공짜로 주지 않는 것이 일반적이다. 위 식은 옵션의 지급액이며, 순이익에는 처음 낸 프리미엄의 만기 가치도 빼야 한다.\par

\par\medskip\noindent\textbf{심화: ‘선형’이라는 표현의 작은 엄밀성}\quad $S_T-K$는 수학적으로 원점을 지나는 선형함수보다는 절편이 있는 아핀함수다. 금융에서는 기초가격에 대한 기울기가 일정하다는 뜻으로 관습적으로 선형 지급액이라고 부른다. 이때 민감도는 $\partial\Pi/\partial S_T=q$다.\par

\subsection{슬라이드 05 · 거래 장소와 계약의 층위}
\textbf{거래소 파생상품}은 수량·만기·결제 단위를 표준화하여 다수 참가자가 호가를 낸다. \textbf{장외(OTC) 파생상품}은 두 당사자가 필요에 맞게 조건을 설계한다. 이는 거래 장소의 구분이다. 한편 ELS·DLF 같은 구조화상품은 투자자가 매수하는 증권·펀드의 지급액 안에 옵션과 채권 등이 \textbf{내재}된 포장 형태다. 따라서 ‘거래소/장외’와 ‘독립 계약/구조화상품’은 서로 다른 분류 축이다.\par

예를 들어 KOSPI200 선물은 거래소에서 표준화되어 거래되고, 특정 기업의 환위험에 정확히 맞춘 환선도는 장외에서 설계할 수 있다. ELS는 여러 지수의 하락 구간에 노출되는 옵션을 사실상 매도한 것과 비슷한 손익을 투자자에게 줄 수 있다. 명칭만으로 원금 보장 여부를 추론하면 안 된다.\par

\subsection{슬라이드 06 · 계약·신용위험·검증가능성}
파생상품은 약속이므로 \textbf{누가 약속을 지키는가}가 중요하다. 장외거래는 양자 신용위험을 직접 관리하고, 청산소를 거치는 거래는 청산소·회원의 보증금과 손실분담 체계가 개입한다. 위험이 사라지는 것이 아니라 담보와 정산 규칙에 따라 다시 배분된다. 계약의 참조 가격과 계산 방식은 독립적으로 검증 가능해야 분쟁이 줄어든다.\par

기초가격에 대한 파생거래의 영향이 무시될 만큼 작다는 ‘외생성’은 분석을 단순하게 하는 가정이다. 대규모 헤지나 만기 직전 포지션은 현물시장에 영향을 줄 수 있어, 현실에서는 이 가정을 확인해야 한다. 또한 다른 계약과 연동된 조기종료 조항이 있으면 하나의 지급액만 따로 계산해서는 안 된다.\par

\subsection{슬라이드 07 · 역사는 위험 이전의 수요가 오래되었음을 보여 준다}
농산물 생산자는 수확 전에 가격을 고정하고 싶고, 구매자는 원재료 비용을 고정하고 싶다. 이런 수요가 오늘날 거래소가 생기기 전부터 선도 형태의 약속을 낳았다. 현대 거래소의 기여는 가격위험을 새로 만든 데 있지 않고, 계약을 표준화하고 매매·청산을 조직해 더 많은 상대방을 만나게 한 데 있다. 고대 문헌의 일화나 최초 거래소의 연대는 ‘현대적 파생계약과 동일했다’는 뜻으로 읽기보다 \textbf{미래 불확실성을 계약으로 다루려는 필요}가 오래되었다는 사례로 읽는 편이 정확하다.\par

\subsection{슬라이드 08 · 세 가지 사용 목적}
\textbf{투기}는 전망에 따라 새로운 위험을 진다. \textbf{차익거래}는 동일한 현금흐름의 가격 차이를 이용한다. \textbf{헤지}는 이미 지닌 위험을 줄인다. 세 목적은 엄격히 분리되지 않는다. 딜러는 고객의 헤지를 받아들여 재헤지하고, 남는 위험을 보유하거나 다른 고객에게 넘긴다. 국내의 ELS와 DLF 사례는 판매 단계의 정보와 실제 손익 구조가 왜 중요한지 보여 준다.\par

\subsection{슬라이드 09 · DLF를 현금흐름으로 분해하기}
금리연계 DLF는 펀드라는 포장 속에 특정 해외 금리·스프레드가 임계값 아래로 내려가면 큰 원금손실이 생기는 구조를 담을 수 있다. 예를 들어 단순화한 만기 지급액을 $N[1+c-\lambda\max(H-X_T,0)]$라 두면, 금리 지표 $X_T$가 $H$보다 낮아질 때 손실이 기울기 $\lambda$로 커진다. 높은 표면 쿠폰 $c$는 이 하방위험을 떠안는 대가다. 이 식은 구조를 설명하는 \textbf{교육용 예시}이며 실제 계약은 관측일·최대손실·중도환매 조건에 따라 달라진다.\par

설명 의무의 경제적 핵심은 판매자와 투자자의 정보 차이다. ‘예상되는 이자’만 보여 주고 손실 구간의 확률·규모와 발행자 위험을 생략하면 투자자는 비교할 수 없다. 국내 금리연계 DLF 논란을 볼 때도 상품 손실과 판매 절차의 문제를 구분해야 한다.\par

식을 숫자로 읽어 보자. 원금 $N=100$만원, 약정 쿠폰 $c=4\%$, 경계 $H=0$, 기울기 $\lambda=10$인 교육용 구조에서 지표가 $X_T=1\%$이면 옵션 손실항은 0이고 지급액은 104만원이다. $X_T=-1\%$이면 $\max(H-X_T,0)=0.01$이므로 지급액은 $100(1.04-10\times0.01)=94$만원이다. 지표가 2\%포인트 더 내려가 $-3\%$가 되면 지급액은 74만원으로 줄어든다. \textbf{지표 1\%포인트 하락이 원금 10\%포인트 손실로 번지는 구간}을 확인해야 표면 쿠폰 4\%의 의미를 알 수 있다. 실제 계약의 지급액에는 원금 하한, 만기·중도환매 조건과 여러 관측일이 포함될 수 있다. \href{https://www.fsc.go.kr/no010101/73971}{금융위원회의 해외금리연계 DLF 점검}.\par

\subsection{슬라이드 10 · 투기는 왜 파생상품으로 하는가}
기초자산을 직접 사지 않아도 작은 초기 담보로 큰 명목 노출을 만들 수 있어, 같은 가격변화가 자기자본 수익률에 크게 반영된다. 예를 들어 100만 원 노출의 선물에 담보 10만 원을 넣고 기초가격이 2\% 떨어지면 계약손실 2만 원은 담보 대비 20\%다. 반대로 2\% 오르면 20\% 수익이지만, 손실은 추가 담보 납입을 요구할 수 있다. \textbf{레버리지의 분모는 계약의 경제적 노출이 아니라 투입 자기자본}이다.\par

전망이 옳아도 현금이 부족하면 만기까지 버티지 못한다. 따라서 방향 예측과 함께 변동성, 손절 기준, 담보 여력, 거래 종료 시점을 정해야 한다.\par

\subsection{슬라이드 11 · 차익거래의 전제와 한계}
같은 미래 지급액 $X_T$를 가진 두 포트폴리오가 오늘 서로 다른 가격이라면 싼 것을 사고 비싼 것을 팔아 차액을 잠글 수 있다. 그러나 ‘비슷한’ 현금흐름만 가진 상대가치 거래는 가격 차이가 더 벌어질 수 있으므로 엄밀한 무위험 차익거래가 아니다. 거래비용, 공매도 제약, 담보 요구, 거래상대방 부도도 실제 실현 가능성을 좌우한다.\par

왜 가격이 같아야 할까? $P_0^{(A)}<P_0^{(B)}$라고 가정하고 오늘 A를 사고 B를 판다면 순현금 $P_0^{(B)}-P_0^{(A)}>0$을 받는다. 만기에는 A에서 $X_T$를 받고 B에 $X_T$를 지급해 어느 상태에서나 순현금 0이다. 오늘 받은 돈을 안전하게 보관할 수 있다면 위험 없이 이익이 확정된다. 따라서 \textbf{동일한 만기, 통화, 지급 조건, 신용위험}을 가진 복제 가능한 현금흐름에는 가격 차이가 지속되기 어렵다. 이 논증에서 A와 B의 만기 지급액이 한 상태라도 다르면 위험이 남는다.\par

\[
X_T^{(A)}=X_T^{(B)}\ \text{모든 상태에서}\quad\Longrightarrow\quad P_0^{(A)}=P_0^{(B)}
\]

이것이 일물일가의 원리다. 뒤에서 선도 가격을 유도할 때 ‘현물을 사서 보유한 포트폴리오’와 ‘선도계약’을 같은 만기 현금흐름으로 만들 것이다.\par

\subsection{슬라이드 12 · 헤지는 전체 손익의 기울기를 바꾼다}
달러 1달러를 만기에 사야 하는 수입업자의 원화 비용은 $S_T$다. 달러 매수 선도의 만기 이익은 $S_T-K$이므로 실제 구매비용에서 이를 빼면 $S_T-(S_T-K)=K$로 고정된다. 이것이 완전헤지다. 절반만 헤지하면 비용은 $\tfrac12S_T+\tfrac12K$로 기울기가 절반 남는다. 비용 급등 위험을 줄이는 동시에 환율 하락의 혜택도 절반 포기한다.\par

일반적으로 필요한 달러 수량을 $A>0$, 선도로 고정한 달러 수량을 $Q\in[0,A]$라고 하면 원화 순비용은 $AS_T-Q(S_T-K)=(A-Q)S_T+QK$다. 환율에 대한 민감도는 $A-Q$이므로 $Q=A$에서만 0이 된다. 예컨대 $A=100$만 달러, $K=1{,}300$원/달러, $Q=50$만 달러면 환율이 $1{,}200$원일 때 비용은 12.5억 원, $1{,}400$원일 때는 13.5억 원이다. 헤지하지 않았다면 각각 12억 원과 14억 원이고, 완전헤지라면 두 경우 모두 13억 원이다. \textbf{변동 폭을 줄이는 대신 좋은 환율의 혜택 일부도 포기한다}는 교환이 숫자로 드러난다.\par

헤지는 ‘이익을 보장’한다기보다 \textbf{원하지 않는 상태 의존성을 줄인다}. 기존 노출의 양과 방향을 쓰지 않고 선도를 거래하면 헤지인지 판단할 수 없다.\par

\subsection{슬라이드 13 · 세 목적을 다시 판별하기}
거래 전에 다음 세 질문을 적자. 기존 현금흐름의 위험이 있는가? 새 계약을 더했을 때 나쁜 상태의 손실이 줄어드는가? 현금흐름을 복제해 오늘 가격 차이를 확정할 수 있는가? 첫째와 둘째가 참이면 헤지, 셋째가 참이면 무차익거래다. 기존 위험이 없는데 새 계약으로 가격 방향에 노출된다면 투기에 가깝다. ELS 투자자가 표면 쿠폰만 보고 하방 옵션 매도의 위험을 보지 못한 경우는 이 구분이 특히 중요하다.\par

\subsection{슬라이드 14 · 다음 질문: 왜 시장이 커졌고 어디서 실패했나}
환율·금리·원자재 가격의 변화가 커지면 기업의 현금흐름을 안정화할 수단에 대한 수요가 커진다. 하지만 계약이 복잡해질수록 거래량의 성장과 위험관리 능력의 성장이 같지는 않다. 다음 장에서는 변동성, 세계화, 전산화가 거래를 늘린 경로와, 현금흐름·담보·조직통제의 실패가 어떻게 대형 손실로 이어졌는지 분리해 본다.\par

\section{03. 변동성, 시장의 성장, 대형 손실}
파생상품 시장이 커진 이유와 손실 사고가 생기는 이유는 서로 연결되어 있다. 가격 변동이 커질수록 위험을 옮길 필요가 커지지만, 계약의 명목금액과 현금 지급 시점이 복잡해지면 장부상 헤지와 실제 현금 부족이 동시에 나타날 수 있다.\par

\subsection{슬라이드 01 · 변동성과 손실을 함께 읽는 법}
이 장에서는 ‘가격위험을 거래하는 시장’이 왜 성장했는지 먼저 본다. 이어서 사례마다 \textbf{원래 노출, 취한 계약, 예상한 상쇄, 실제 실패 지점}을 같은 순서로 점검한다. 손실액만 외우면 왜 같은 일이 반복되는지 이해하기 어렵다.\par

\subsection{슬라이드 02 · 사례를 보는 네 가지 렌즈}
사고는 크게 가격위험을 잘못 잰 경우, 만기·기초자산이 맞지 않는 헤지, 증거금 때문에 유동성이 먼저 고갈된 경우, 조직 안의 승인·보고가 무너진 경우, 판매자가 복잡성을 이용한 경우로 나눌 수 있다. 한 사건에는 보통 여러 원인이 겹친다. 예컨대 선물을 사용한 헤지가 최종 만기 손익을 줄여도 \textbf{중간 현금흐름}의 위기는 남는다.\par

\subsection{슬라이드 03 · 변동성은 왜 위험 이전 수요를 만든다}
기업의 이익을 $\Pi=P Q-C$라고 하면 판매가격 $P$, 수량 $Q$, 원재료 비용 $C$의 불확실성이 이익 변동을 만든다. 환율·금리·원자재 가격이 변하면 기존 사업의 현금흐름에도 영향을 준다. 선도와 스왑은 한 주체의 변동 노출을 다른 주체에게 이전한다. 변동성 증가 자체가 항상 거래량 증가를 뜻하지는 않지만, 위험을 고정하려는 의사와 위험을 받아 프리미엄을 얻으려는 의사가 동시에 있을 때 거래가 늘어난다.\par

\subsection{슬라이드 04 · 고정환율에서 변동환율로}
고정환율제에서는 특정 구간에서 환율의 일일 변동이 억제될 수 있다. 변동환율제가 확대되면 기업이 미래 외화 매출을 원화로 환산할 때 받을 금액의 불확실성이 눈에 띈다. 이때 3개월 후 달러를 받을 수출업자는 달러 매도 선도로 원화 수입을 고정할 수 있다. 환율 변동 그래프와 파생상품의 등장을 나란히 놓은 그림은 \textbf{위험의 등장 → 계약 수요}라는 경제적 동기를 보여 준다. 다만 상관관계만으로 특정 계약이 변동성을 낮췄다고 단정할 수는 없다.\par

\subsection{슬라이드 05 · 환율 그래프의 세로축을 읽는 방법}
여러 통화의 환율 변화율을 한 그래프에 그릴 때 세로축은 보통 가격 수준이 아니라 기간별 수익률이다. 원/달러가 1,200원에서 1,260원으로 오르면 달러의 원화가격 수익률은 $1,260/1,200-1=5\%$이다. 그러나 원화의 달러가치는 반대로 내려간다. \textbf{표시 통화의 방향}을 확인하지 않으면 상승·하락의 경제적 의미를 거꾸로 읽는다. 그래프의 큰 관측치는 평균보다 더 중요한 담보·손절 문제를 일으킬 수 있다.\par

\subsection{슬라이드 06 · 금리 변동은 채권의 가격위험이다}
금리가 오르면 고정 쿠폰 채권의 현재가치는 떨어진다. 단일 확정 현금흐름 $C$의 가격은 $P=C e^{-rT}$이고, 작은 금리변화에 대해 $\Delta P\approx -T P\Delta r$다. 이 근사식의 $T$가 긴 채권일수록 같은 금리 변화에 가격이 더 민감하다는 직관을 준다. 실제 채권은 여러 현금흐름의 가중평균 만기인 듀레이션으로 일반화한다.\par

\subsection{슬라이드 07 · 금리 그래프는 ‘수준’과 ‘변화’를 구별한다}
10년물 금리의 변화 $\Delta y_t=y_t-y_{t-1}$를 그린 그래프가 위아래로 흔들린다는 것은 금리 수준이 음수라는 뜻이 아니다. 변화가 커질수록 채권 가격과 차입비용의 예상 오차가 커진다. 장기 고정금리 채권을 가진 기관은 금리선물 매도나 고정금리 지급·변동금리 수취 스왑으로 금리 상승 시 손실을 일부 상쇄할 수 있다. 헤지 수량은 명목원금만 맞춰서는 안 되고 금리 민감도도 맞춰야 한다.\par

\subsection{슬라이드 08 · 원자재의 가격위험은 생산과 연결된다}
항공사는 항공유 가격 상승, 정유사는 원유와 정제품 간 가격 차이, 식품회사는 곡물 가격 상승에 노출된다. 원자재 선물은 미래 구매·판매 가격을 일부 고정하지만 품질·지역·인도월이 실제 현물거래와 다르면 베이시스 위험이 남는다. 금속이나 석유는 보관비용과 재고가 부족할 때 즉시 사용할 수 있는 이익인 편익수익률도 가격에 영향을 준다.\par

\subsection{슬라이드 09 · 원유와 금의 변동을 같은 숫자로 비교할 때}
같은 기간의 가격 변화율을 그리면 두 자산의 충격 크기를 비교할 수 있다. 그러나 한 자산의 표준편차가 크다는 사실만으로 헤지가 더 쉽거나 어렵다고 결론 내릴 수 없다. 중요한 것은 기업의 실제 현금흐름과 사용할 선물가격의 \textbf{공분산}이다. 원유 가격이 크게 움직여도 사업 비용과 거의 연동되지 않으면 그 선물은 좋은 헤지가 아니다.\par

\subsection{슬라이드 10 · 세계화·전산화·거래비용의 결합}
매출은 달러, 비용은 원화, 부채는 엔화로 발생하는 기업은 하나의 환율만 고정해도 다른 환율 위험이 남는다. 세계화는 노출의 종류를 늘렸고 전산화는 복잡한 계약의 가격과 위험을 빠르게 계산·거래하게 했다. 표준화와 낮아진 거래비용은 작은 가격 차이에도 거래를 가능하게 한다. 반대로 모형이 숫자를 빠르게 만든다는 사실이 모형의 전제나 극단 상황까지 옳다는 뜻은 아니다.\par

\subsection{슬라이드 11 · 명목원금과 시장가치는 다른 숫자다}
스왑 명목원금 $N$은 이자를 계산할 기준금액이다. 1억 원 명목의 금리스왑이 있다고 1억 원이 처음 오가거나 전액이 손실 위험에 놓이는 것은 아니다. \textbf{총시장가치}는 이미 체결한 계약을 현재 시장조건으로 대체하는 데 드는 양(+)·음(-)의 가치의 절댓값을 합한 지표다. 둘 다 시장 규모를 보는 데 쓰지만 신용위험을 단독으로 정확히 나타내지 않는다. 법적 상계와 담보를 고려한 노출은 또 다르다.\par

\[
\text{명목원금}=N,\qquad \text{총시장가치}=\sum_i |V_i|,\qquad
\text{순신용노출}\ne\sum_i |V_i|.
\]

표의 1999·2001·2010년 수치는 당시 시장의 역사적 스냅샷이다. 서로 다른 조사 범위와 집계 방식, 중앙청산 시 계약 분할 효과를 확인한 뒤 기간별 증가율을 해석해야 한다.\par

\subsection{슬라이드 12 · 금리 파생상품이 큰 비중을 차지하는 이유}
금리선도계약(FRA), 금리스왑, 옵션은 채권과 대출의 금리위험을 교환한다. 금리스왑은 명목원금을 교환하지 않고 여러 날짜의 고정·변동 이자 차액을 지급하는 경우가 많다. 은행의 장기 대출과 단기 예금처럼 만기가 다른 자산·부채는 금리 변화에 민감하므로 거래 수요가 크다. 표에서 명목원금이 큰 항목을 \textbf{최대 손실액}으로 읽지 말고, 어느 현금흐름의 기준액인지 먼저 보자.\par

\subsection{슬라이드 13 · 거래소 계약 수와 장외 명목원금은 비교 단위가 다르다}
거래소 그래프의 계약 수는 거래 횟수 또는 미결제 계약 수와 관련 있고, 장외 통계의 명목원금은 계약별 기준액의 합이다. 계약 1개의 규모가 상품마다 달라 두 숫자를 바로 비교할 수 없다. 시장의 깊이를 알고 싶다면 거래량·미결제약정·호가 스프레드·호가 잔량을 함께 봐야 한다. 세계 여러 거래소의 성장도 위험 분산의 기회와 시장 간 전염 통로를 동시에 늘린다.\par

\subsection{슬라이드 14 · 비유가 놓치는 것: 손실의 경로}
큰 손실은 파생상품을 사용했다는 사실만으로 생기지 않는다. 작은 담보로 큰 노출을 가진 상태에서 불리한 가격 움직임이 생기고, 일일정산이나 담보 재요구가 들어오며, 자산을 급히 팔아야 하는 \textbf{경로}에서 위기가 커진다. 만기의 예상 이익이 양(+)이어도 오늘 납입할 현금이 없으면 거래를 유지하지 못한다.\par

\subsection{슬라이드 15 · 파생상품과 실패 원인을 분리하기}
차입해서 장기채권을 사는 전략은 파생계약 없이도 금리 상승과 단기자금 조달 실패에 취약하다. 반대로 파생상품은 그 위험을 헤지할 수도 있다. 사건을 볼 때 계약의 종류만으로 책임을 돌리기보다 (1) 누가 어떤 위험을 보유했는지, (2) 그 위험이 계약으로 줄었는지, (3) 조직이 실제 규모를 알고 있었는지 묻는다. \textbf{상품 구조, 레버리지, 유동성, 내부통제}는 각각 별개의 원인이다.\par

\subsection{슬라이드 16 · 1990년대 손실 목록: 서로 다른 고리}
스왑 손실인 Gibson Greetings와 Procter \& Gamble에서는 구조화된 금리 지급액과 가격평가가 핵심이었다. Metallgesellschaft의 에너지 선물에서는 장기 고객 약정과 단기 선물 헤지 사이의 만기 불일치가 중요했다. Codelco의 구리 선물이나 여러 환옵션 사례는 기초가격에 대한 방향성 노출을 보여 준다. 같은 ‘파생상품 손실’ 표의 한 줄이라도 실패 고리가 달라 공통 원인을 한 단어로 묶을 수 없다.\par

\subsection{슬라이드 17 · 손실 목록: 레버리지와 통제}
Orange County는 차입 확대와 금리상승에 대한 취약성이, Barings는 승인되지 않은 거래와 손실 은폐가 대표적이다. 선물의 일일정산은 손실을 조기에 드러내지만, 내부 보고가 가려지면 청산소에서 보이는 담보 요청을 경영진이 위험 경고로 활용하지 못한다. 표의 일부 항목에는 파생상품 외에도 레포와 구조화채권이 함께 들어 있으므로 \textbf{계약 유형}과 \textbf{손실 메커니즘}을 별도로 분류해야 한다.\par

\subsection{슬라이드 18 · 손실 목록: 수렴 거래의 함정}
LTCM의 상대가치 포지션은 비슷한 자산의 가격 차이가 다시 좁혀질 것으로 기대했다. 그러나 비슷한 현금흐름은 동일한 현금흐름이 아니다. 위기 때 위험회피가 커지면 차이가 더 벌어지고, 담보 추가 납입 때문에 포지션을 줄이는 매매가 다시 가격 차이를 키운다. 확률상 드문 손실이 \textbf{레버리지와 유동성}을 통해 생존 문제로 바뀌는 과정이다.\par

\subsection{슬라이드 19 · 사고를 설명하는 위험의 분류표}
\textbf{시장위험}은 기초가격 자체의 변화, \textbf{베이시스 위험}은 헤지 대상과 헤지 수단의 가격 차이, \textbf{자금유동성 위험}은 담보나 채무에 쓸 현금 부족, \textbf{시장유동성 위험}은 큰 가격충격 없이 포지션을 팔기 어려움이다. \textbf{모형위험}은 분포·상관관계·평가식이 잘못된 경우, \textbf{운영위험}은 권한·보고·감시가 실패한 경우다. 판매 단계에서는 제조자·판매자·투자자 간 정보비대칭이 추가된다.\par

\par\medskip\noindent\textbf{심화: 자금유동성과 시장유동성이 서로를 증폭시키는 과정}\quad 증거금 납입을 위해 자산을 매도하면 매도 압력으로 가격이 내려간다. 가격 하락은 다른 참가자의 평가손실과 담보 요청을 늘려 또 다른 매도를 유발한다. 개별 계약의 만기 손익을 계산하는 것만으로 이 피드백을 알 수 없으므로 스트레스 시나리오에 현금조달 계획을 함께 넣어야 한다.\par

\subsection{슬라이드 20 · Metallgesellschaft: 경제적 헤지와 현금 부족}
미국 자회사는 고객에게 여러 해 동안 석유제품을 고정가격에 공급하기로 했으므로 현물 가격 상승에 불리한 \textbf{장기 매도 약정}을 보유했다. 가격 상승을 상쇄하려고 가까운 만기의 원유 선물을 매수해 만기마다 다음 선물로 옮기는 \texttt{stack and roll} 방식을 썼다. 유가가 하락하자 장기 고객 약정에는 장래의 경제적 이익이 생겼지만, 매수 선물에서는 \textbf{즉시 현금으로 낼 손실}이 발생했다. 두 현금흐름의 시점과 담보 요구가 달라 헤지가 자금위기를 막지 못했다.\par

이 전략은 고객 계약과 선물이 같은 기초품질·만기·수량이 아니므로 완전 복제도 아니다. 단기 선물을 장기 노출에 계속 굴릴 때 다음 계약 가격이 어떻게 바뀔지 알 수 없는 롤오버 위험이 남는다.\par

현금 시점의 차이를 작은 예로 확인하자. 고객에게 $T_2$에 100배럴을 배럴당 $80$에 공급하기로 하고, 가까운 만기 $T_1$의 선물을 100배럴 매수했다고 하자. $T_1$에 가까운 만기와 장기 인도 가격이 모두 $80$에서 $60$으로 내려간 단순한 상황이라면, 선물의 가격변화 손익 약 $100(60-80)=-\$2{,}000$은 일일정산을 통해 먼저 현금으로 나간다. 반면 고객 약정에서 저렴해진 기초물을 사서 공급할 수 있는 이익은 $T_2$의 이행 또는 계약 재평가에 걸쳐 나타난다. \textbf{경제적 가치의 상쇄와 같은 날 사용할 수 있는 현금의 상쇄는 다르다.} 실제로는 두 만기의 가격변화도 정확히 같지 않아 베이시스·롤오버 위험이 더해진다. \href{https://papers.ssrn.com/sol3/papers.cfm?abstract_id=6104}{헤지 만기구조 분석}.\par

\subsection{슬라이드 21 · Gibson Greetings: 제곱 금리 지급액}
기업은 처음에 고정금리 부채를 변동금리로 바꾸려는 합리적인 동기를 가질 수 있다. 그러나 이후 들어간 구조화 스왑은 지급률이 대략 $L^2/0.06$처럼 시장금리 $L$의 제곱에 반응했다. $L=6\%$일 때 이 항은 $6\%$지만, $L=9\%$면 $13.5\%$가 되어 3\%포인트 금리 상승이 7.5\%포인트의 지급률 증가를 만든다. 더 높은 현재 수취금리와 맞바꾼 \textbf{비선형 금리상승 위험}이다.\par

손실을 확정하기 싫어 더 복잡한 거래로 미루면 기초 부채와의 연결이 약해진다. 미국 SEC는 당시 판매자가 손실을 과소평가한 평가값을 제시한 문제를 지적했다. 독립 평가와 거래 목적의 재확인이 필수인 이유다. \href{https://www.sec.gov/newsroom/speeches-statements/tsty1598-testimony-otc-derivatives-us-financial-markets}{SEC의 사건 설명}.\par

\subsection{슬라이드 22 · 평가자료가 읽기 어려울 때 확인할 것}
구조화 스왑의 제안서나 보도 문구가 촘촘할수록 \textbf{조건부 지급액을 식과 시나리오로 다시 적는 작업}이 중요하다. 금리가 $5\%,6\%,9\%$일 때 각 기간에 누가 얼마를 지급하는지 표로 만들고, 일찍 종료할 때의 현재가치를 따로 계산한다. 제안서의 ‘절감되는 이자’가 초기 기간에만 적용될 수 있으며, 중도해지 금액이 미래 불리한 현금흐름의 현재가치를 이미 반영할 수 있다. 판매자의 평가값 하나만으로 위험을 판단하지 않는다.\par

\subsection{슬라이드 23 · Orange County: 역레포와 역변동금리채}
지자체 투자풀은 역레포로 보유 채권을 담보 삼아 단기자금을 빌리고 더 많은 중장기 채권을 샀다. 금리가 오르면 채권 가격이 하락하고, 차입 비용도 오르며, 담보가치가 떨어져 추가 담보 요구가 생긴다. 역변동금리채는 금리가 오를수록 쿠폰이 줄어들어 가치가 더욱 민감하게 하락할 수 있다. 이는 파생상품만의 문제가 아니라 \textbf{장기자산을 단기차입으로 레버리지한 구조}의 문제다.\par

\[
\frac{\Delta P}{P}\approx -D_{\rm mod}\Delta y
\]

여기서 수정듀레이션 $D_{\rm mod}$는 금리 변화에 대한 가격의 1차 민감도다. 레버리지가 커지면 같은 채권가격 하락이 자기자본에 더 크게 반영된다.\par

역변동금리채의 쿠폰을 단순화해 $c_t=a-bL_t$($b>0$)라고 쓰면, 시장금리 $L_t$가 오를 때 \textbf{받는 쿠폰은 줄고 할인율은 올라간다.} 채권가치에는 두 방향 모두 불리하다. 여기에 자기자본 1억 달러와 단기차입 1억 달러로 2억 달러의 채권을 보유한 교육용 사례를 생각해 보자. 자산가치가 10\% 떨어져 1.8억 달러가 되면 차입금 1억 달러를 뺀 자기자본은 0.8억 달러로 20\% 감소한다. 담보권자가 추가 담보를 요구하면 만기까지 기다려 회복될 가능성과 별도로 오늘 자산을 팔아야 할 수 있다. 당시 금리 상승이 차입비용·역변동금리채 이자·채권가격·담보 요구를 동시에 악화시켰다는 점은 \href{https://www.sec.gov/enforcement-litigation/administrative-proceedings/33-8141}{미국 SEC의 사건 기록}에서 확인할 수 있다.\par

\subsection{슬라이드 24 · Sumitomo: 집중과 권한의 위험}
구리 현물·선물·옵션·스왑을 한 거래자에게 집중시키고 브로커 계좌 개설과 자금 이동까지 폭넓은 권한을 주면 독립적인 견제가 어렵다. 거대한 포지션은 자기 매매가 시장가격에 영향을 미쳐 장부 평가와 청산 가능 가격이 달라지는 \textbf{시장충격 위험}도 만든다. 투자 판단의 옳고 그름 외에 포지션 한도, 브로커 확인, 거래와 결제의 직무 분리가 중요한 사례다.\par

\subsection{슬라이드 25 · Barings: 현물 차익거래처럼 보인 방향성 베팅}
서로 다른 거래소의 닛케이 선물 가격 차이를 이용하는 거래라면 한쪽 매수와 다른 쪽 매도가 함께 있어야 한다. 실제로는 손실을 숨긴 채 주가지수 상승에 베팅하는 포지션이 커졌다. 거래자가 거래와 사후처리를 함께 통제하고 손실계좌를 숨긴 것이 내부통제 실패의 핵심이다. 단순히 ‘선물은 위험하다’는 결론보다 \textbf{보고된 전략과 실제 순노출의 차이를 확인해야 한다}는 교훈이 더 구체적이다.\par

정상적인 거래소 간 스프레드라면 한 시장의 선물 100계약 매수와 다른 시장의 유사 선물 100계약 매도가 함께 기록되어 지수 전체의 상승·하락에 대한 민감도가 작아져야 한다. 매도 기록이 없거나 허위이면 장부상의 ‘상대가격 거래’는 사실상 100계약 순매수다. 따라서 계약별 상대방 확인과 증거금 송금 내역을 독립 부서가 대조해야 한다. 한 사람이 거래·결제·오류계좌를 모두 통제하면 손익표만으로 이 차이를 발견하기 어렵다. \href{https://www.gov.uk/government/publications/report-into-the-collapse-of-barings-bank}{영국 은행감독위원회의 조사 보고서}.\par

\subsection{슬라이드 26 · Barings: 스트래들 매도의 꼬리위험}
콜과 풋을 함께 매도하는 스트래들은 기초가격이 행사가격 근처에 머물면 프리미엄을 얻지만, 크게 움직이면 한쪽 옵션의 손실이 커진다. 닛케이 급락은 주가지수 선물 매수와 스트래들 매도에 모두 불리했다. 손실을 만회하려 포지션을 늘리면 나쁜 상태에서 노출이 더 커진다. 선물의 증거금 요청은 경고 신호였으나, 본사 자금 지원이 실제 포지션 검증 없이 이어지면 위험을 감춘 채 거래를 연장하게 된다.\par

\[
\Pi_T^{\rm short\ straddle}=C_0+P_0-\max(S_T-K,0)-\max(K-S_T,0)
=C_0+P_0-|S_T-K|.
\]

이 식에서 $C_0+P_0$는 처음 받은 프리미엄이고, $|S_T-K|$는 만기에 콜과 풋 중 가치가 있는 쪽의 지급액이다. 금리와 거래비용을 잠시 무시하면 손익분기점은 $S_T=K\pm(C_0+P_0)$이며, 그 범위를 벗어나면 손실이 커진다. 지수가 급락하면 \textbf{주가지수선물 매수 손실}과 \textbf{매도한 풋옵션 손실}이 같은 방향으로 누적된다. 여러 계약을 묶었다는 사실만으로 위험이 분산되는 것은 아니며, 나쁜 상태에서 각 계약의 손익 부호가 같은지 확인해야 한다. \href{https://www.gov.uk/government/publications/report-into-the-collapse-of-barings-bank}{영국 은행감독위원회의 조사 보고서}.\par

\subsection{슬라이드 27 · LTCM: 수렴 거래는 확정 차익이 아니다}
수익률이 비슷한 채권 두 개를 사고팔아 스프레드 축소를 기대하면 정상 시기에 작은 수익이 반복될 수 있다. 작은 스프레드에 의미 있는 수익률을 내기 위해 레버리지를 쓰면 반대 방향의 큰 변동에는 취약하다. 위기 때 ‘안전자산 선호’로 스프레드가 벌어지면 포지션의 손익이 나빠진다. 같은 만기 현금흐름을 완벽히 복제한 거래가 아니므로 손실이 일시적이라고 단정할 수 없다.\par

\subsection{슬라이드 28 · LTCM: 증거금 요구가 가격을 다시 움직인다}
포지션이 불리해지면 거래상대방이 담보를 더 요구한다. 현금이 부족해 포지션을 팔면 유동성이 얕은 시장의 가격이 더 나빠지고, 남은 포지션도 추가 손실을 입는다. 이는 개별 모형의 상관계수만으로는 충분히 포착하기 어렵다. 1998년 뉴욕 연준은 시장의 급격한 청산을 우려해 민간 채권자들의 협의를 촉진했다. \textbf{연준이 자기 자금으로 손실을 메운 것은 아니다.} \href{https://www.federalreservehistory.org/essays/ltcm-near-failure}{연방준비제도 역사 자료}.\par

\subsection{슬라이드 29 · LTCM: 모형위험과 민간 구조조정}
14개 금융회사가 약 36억 달러를 투입해 지분을 취득하고 포지션을 질서 있게 줄이는 구조가 마련되었다. 모형의 과거 데이터가 위기 때의 동시 매도와 거래비용을 충분히 반영하지 못하면 VaR 수치가 안전을 보증하지 못한다. 상대방마다 제공한 신용한도가 합쳐져 한 펀드의 시스템 전체 노출이 커질 수도 있다. \href{https://www.federalreservehistory.org/essays/ltcm-near-failure}{연방준비제도 역사 자료}.\par

\subsection{슬라이드 30 · Société Générale: 존재하지 않는 헤지}
지수선물 매수와 매도를 함께 가진 차익거래라고 보고했지만 매도 쪽이 허위라면 실제 포지션은 거대한 단방향 매수다. 위험관리 시스템이 상계 후 순노출만 확인하고 각 거래의 실재성과 상대방 확인을 놓치면 허위 헤지로 한도를 회피할 수 있다. 따라서 시스템의 숫자뿐 아니라 \textbf{독립적인 거래 확인, 취소 기록, 실제 증거금 흐름}을 대조해야 한다.\par

예컨대 실제 선물 매수 명목액이 50억 유로이고 시스템에만 같은 규모의 가짜 매도가 들어 있다면, 보고된 순노출은 0이지만 경제적 노출은 50억 유로 매수다. 지수가 1\% 떨어질 때 실제 가격손실은 단순 근사로 5천만 유로다. 이 숫자는 사건의 실제 포지션 규모가 아니라 \textbf{상계 장부가 왜 위험을 숨길 수 있는지} 보여 주는 예시다. 사후 통제는 총매수와 총매도를 각각 확인하고, 매도 거래의 상대방·청산 기록이 존재하는지 독립적으로 대조해야 한다. \href{https://www.societegenerale.com/en/news/newsroom/explanatory-note-about-exceptional-fraud}{Société Générale의 사건 설명}.\par

\subsection{슬라이드 31 · ELS: 쿠폰 뒤에 있는 하방위험}
주가연계증권(ELS)은 발행자의 신용위험을 지닌 증권에 지수·주가 옵션의 지급 구조가 결합된다. 여러 지수 중 가장 성과가 나쁜 지수에 원금상환이 연동되는 \texttt{worst-of} 구조에서는 각 지수가 독립적으로 조금만 떨어져도 가장 낮은 지수가 손실을 결정할 수 있다. 녹인 배리어를 터치했는지, 조기상환 조건을 충족했는지에 따라 지급액이 달라지는 \textbf{경로 의존성}도 중요하다. 표면 쿠폰은 이런 손실 위험의 대가이며 예금 이자와 다르다.\par

예를 들어 두 지수 중 하나라도 최초 수준의 50\% 미만을 기록하고 만기에 최악 지수가 70\%면, 단순한 교육용 구조에서는 약 30\%의 원금손실이 날 수 있다. 실제 손익은 계약별 조건을 읽어야 한다. 홍콩 H지수 연계 ELS는 지수 하락과 조기상환 지연이 왜 중요한지 보여 주는 국내 사례다. \href{https://www.fsc.go.kr/po010106/81314}{금융위원회 H지수 ELS 현황}.\par

두 지수의 만기 성과를 $R_1=S_{1,T}/S_{1,0}$, $R_2=S_{2,T}/S_{2,0}$로 두면 최악 지수 성과는 $R_{\min}=\min(R_1,R_2)$다. 만기 성과가 경계 $h$ 아래로 \textbf{하나라도} 떨어질 확률은 $P(R_{\min}<h)=1-P(R_1\ge h, R_2\ge h)$다. 두 지수의 경계 하락 확률이 각각 10\%이고 독립이라는 교육용 가정에서는 $1-0.9^2=19\%$로, 지수 하나의 10\%보다 높다. 실제 지수들은 독립이 아니므로 19\%를 실제 손실확률로 사용해서는 안 된다. 녹인 여부가 중간 경로에 달린 ELS라면 만기 확률만으로 지급액을 구할 수도 없다. 이 계산은 높은 쿠폰을 볼 때 ‘둘 중 더 나쁜 지수’와 \textbf{지수 간 상관관계}를 함께 살펴야 하는 이유를 보여 준다.\par

\subsection{슬라이드 32 · DLF 손익 그림의 꺾이는 구간}
금리나 신용지표가 정한 구간 위에 있으면 쿠폰이 지급되다가 기준 이하로 내려가면 지급액이 급격히 줄 수 있다. 그림을 읽을 때 가로축이 금리 \textbf{수준}인지 변화폭인지, 세로축이 원금상환율인지 이자율인지 확인한다. ‘약정 이자’는 모든 상태에서 지급되는 확정수익이 아니라 특정 조건을 만족한 상태의 지급액이다. 손익선의 기울기가 갑자기 달라지는 경계에서 작은 지표 변화가 큰 손실 차이를 만들 수 있다.\par

\subsection{슬라이드 33 · 잘못 판매된 상품의 정보 구조}
제조자는 옵션의 지급액과 헤지 비용을 알고, 판매자는 고객의 목적·지식·손실 감내 능력을 확인하며, 투자자는 최종 위험을 떠안는다. 각 단계에서 정보가 끊기면 표면 금리만 전달되고 \textbf{손실 조건·확률·발행자 신용위험·중도매도 비용}이 누락된다. 정보비대칭을 줄이는 실무적인 방법은 상품명을 설명하는 것보다 극단·중간·기준 시나리오의 원금상환액을 숫자로 보여 주는 것이다.\par

\subsection{슬라이드 34 · 양날의 검을 관리 기준으로 바꾸기}
파생상품은 위험을 줄이거나 늘릴 수 있다. 관리자는 명목원금뿐 아니라 순델타, 금리민감도, 스트레스 손실, 추가 담보 요구, 만기별 현금유출을 같이 보아야 한다. 현금흐름이 복잡한 상품일수록 모형이 수익률의 평균을 잘 맞춘다는 사실보다 손실 꼬리와 조기 종료 조건을 잘 설명하는지가 중요하다. 규제와 청산제도는 이 위험을 줄이지만 모든 계약을 무위험으로 만들지는 않는다.\par

\subsection{슬라이드 35 · 작은 수익과 큰 손실의 비대칭}
옵션 매도처럼 대부분의 기간에 작은 프리미엄을 얻고 드문 큰 가격변화에서 손실을 입는 전략은 과거 평균 수익률만으로 판단하면 위험해 보이지 않을 수 있다. 따라서 ‘한 달 1\% 수익’ 같은 요약보다, 특정 위기 상태에서 \textbf{최대 자금 필요액}과 강제 청산 여부를 먼저 계산해야 한다. 작은 이익을 반복하는 포지션의 핵심 검사는 꼬리 사건의 손실 규모다.\par

\subsection{슬라이드 36 · 선도와 선물로 돌아가기}
앞선 사례들은 기본 계약을 정확히 알아야 복잡한 계약도 이해할 수 있음을 보여 준다. 다음 장에서는 선도가 만기에 한 번 정산되는 것과 선물이 매일 정산되는 것의 차이를 거래 구조·증거금 표로 확인한다. 이후 무차익 가격을 유도하고, 마지막으로 기업의 기존 현금흐름에 선도를 더한 헤지를 수학적으로 설계한다.\par

\section{04. 선도·선물의 계약과 결제}
선도와 선물의 만기 손익은 비슷한 직선으로 나타낼 수 있다. 중요한 차이는 그 손익을 \textbf{언제}, \textbf{누구를 통해}, \textbf{어떤 담보 체계 아래} 현금으로 주고받느냐에 있다. 여기서는 계약 규격에서 청산소·일일정산·증거금 호출까지 순서대로 연결한다.\par

\subsection{슬라이드 01 · 선형 계약의 출발점}
기초자산 한 단위를 만기 $T$에 가격 $K$로 사는 선도의 매수자는 $S_T-K$, 매도자는 $K-S_T$의 만기 이익을 얻는다. 선물은 거래소가 정한 규격으로 비슷한 가격위험을 거래하지만 매일 손익을 정산한다. 따라서 만기의 선 하나만 그리면 중간 현금 필요액이라는 중요한 차이가 보이지 않는다.\par

\subsection{슬라이드 02 · 거래의 전 과정}
선물거래는 \textbf{상품 규격 확인 → 주문·체결 → 청산소를 통한 포지션 등록 → 초기증거금 납입 → 일일정산 → 반대매매 또는 만기결제}로 진행된다. 각 단계는 다른 질문에 답한다. 상품 규격은 ‘무엇을 얼마나’에, 청산은 ‘누가 이행을 보증하는가’에, 증거금은 ‘불리한 가격 변화에 현금을 어떻게 확보하는가’에 답한다.\par

\subsection{슬라이드 03 · 선물 계약의 규격}
선물이 대량 거래되려면 기초물의 품질, 계약 크기, 가격 표시 단위, 호가단위, 인도 가능한 장소와 기간이 미리 정해져야 한다. 옥수수 선물 1계약이 몇 부셸인지 알아야 호가가 $0.01$ 움직일 때 손익을 계산할 수 있다. 일반식은 $\Delta\Pi=q\,m\,\Delta F$다. 여기서 $q$는 계약 수, $m$은 계약당 기초수량, $\Delta F$는 가격변화다. 품질 대체가 허용되면 인도 가격 조정도 계약 규격에 포함된다.\par

\subsection{슬라이드 04 · 고객에서 청산소까지}
일반 고객은 선물중개회사(FCM)를 통해 주문하고, 중개회사는 거래소 회원과 청산회원을 통해 계약을 청산한다. 거래소는 거래 규칙과 매매 시장을 제공하고, 청산소는 체결 뒤 계약의 이행 관계를 관리한다. 고객이 누구와 화면상 거래했는지와 \textbf{신용상 누구에게 이행을 요구하는지}가 달라진다. 청산소가 개입해도 고객은 중개회사, 청산회원, 청산소의 규칙을 차례로 따른다.\par

\subsection{슬라이드 05 · 청산소의 법적·경제적 기능}
청산소가 없으면 매수자와 매도자가 서로의 만기 이행을 직접 걱정한다. 청산소가 매수자에게는 매도자, 매도자에게는 매수자 역할을 하도록 계약을 대체하면 직접 양자 신용위험을 줄이고 다수 거래를 효율적으로 상계할 수 있다. 다만 청산소의 약속은 \textbf{초기·변동 증거금, 회원 자본, 손실분담 절차}로 지탱된다. ‘청산소가 있으므로 부도 가능성이 0’이라고 해석하면 안 된다. \href{https://www.cftc.gov/LearnAndProtect/EducationCenter/economicpurpose}{CFTC의 청산·증거금 설명}.\par

\subsection{슬라이드 06 · 기록의 두 층: 고객 계좌와 회원 계좌}
고객 A가 100계약 매수, 고객 B가 90계약 매도해도 중개회사는 두 고객의 권리를 각각 기록한다. 청산소는 회원별로 긴 포지션과 짧은 포지션을 관리하고, 허용되는 범위에서 노출을 상계한다. \textbf{총거래량}은 그날 체결된 계약 수, \textbf{미결제약정}은 남아 있는 매수·매도 짝의 수다. 매수 120계약과 매도 120계약이 남아 있으면 미결제약정은 120계약이지 240계약이 아니다.\par

같은 회원 안에서 장단기 포지션을 상계해 요구 증거금이 줄 수 있어도 각 고객의 손익과 채무는 사라지지 않는다. 집계표를 볼 때 고객 단위와 회원 단위를 혼동하지 말자.\par

\subsection{슬라이드 07 · 시장 참가자의 역할}
중개인은 고객 주문을 집행하고, 자기계정 거래자는 스스로 위험을 진다. 단기 매매자와 장기 포지션 보유자의 차이는 \textbf{보유 기간}이지 법적 계약의 종류가 아니다. 상품거래자문업자(CTA)는 투자 판단을 자문하고, 공동 투자기구의 운영자는 투자자 돈을 모아 운용한다. 참가자마다 원하는 위험 방향과 시간범위가 다르므로 시장에 양쪽 호가가 생긴다.\par

\subsection{슬라이드 08 · 초기증거금·유지증거금·일일정산}
초기증거금 $IM$은 계약을 열 때 요구되는 이행 담보다. 유지증거금 $MM$ 아래로 계좌가 내려가면 추가 납입을 요구한다. \textbf{증거금은 기초자산의 매입대금이나 선물의 옵션 프리미엄이 아니다.} 매일 정산가격 $F_t$가 정해질 때, 매수 1계약의 일일손익은 $m(F_t-F_{t-1})$이고 그만큼 계좌가 늘거나 줄어든다. 가격 제한폭과 포지션 한도는 급격한 시장 변화와 과도한 집중을 관리하지만, 손실을 없애지 않는다. \href{https://www.cftc.gov/LearnAndProtect/EducationCenter/economicpurpose}{CFTC 설명}.\par

예를 들어 계약당 기초수량이 $m=5{,}000$인 선물 2계약을 매수하고 정산가격이 하루에 $0.03$ 하락하면 일일손익은 $2\times5{,}000\times(-0.03)=-\$300$이다. 기초자산 1만 단위를 실제로 산 것이 아니라 \textbf{그 수량의 가격변화에 해당하는 금액}을 증거금 계좌에서 지급한 것이다. 반대로 $0.03$ 상승하면 같은 금액이 계좌에 들어온다. 가격변화를 다음 날로 넘기지 않고 오늘 정산한다는 점이 선물의 자금수요를 만든다.\par

증거금 계좌의 잔액과 투자자의 순손익도 분리해야 한다. 추가로 넣은 $600$는 계좌 사이의 자금 이동이며, 그 자체가 새로운 손실은 아니다. 계약을 닫을 때 남은 담보를 돌려받고, 그동안의 일일정산액과 담보 조달비용을 합쳐 최종 손익을 계산한다. 따라서 같은 만기 지급액을 가진 선도와 선물이라도 선물을 유지할 \textbf{중간 현금 여력}이 다를 수 있다.\par

\subsection{슬라이드 09 · 옥수수 선물의 계약금액과 담보}
4계약 합계 20,000부셸을 부셸당 $2.21$에 매수하면 명목 계약금액은 $44,200$이다. 초기증거금이 $1,892$라면 명목금액/담보금은 약 $23.36$이다. 이것은 가격이 1\% 움직일 때 담보의 약 23\%에 해당하는 손익이 생길 수 있다는 단순 레버리지 지표다. 실제 위험은 변동폭과 추가 납입 가능액에 달려 있다. 가격이 $0.02$ 하락하면 일일손실은 $20,000\times(-0.02)=-\$400$이다.\par

\subsection{슬라이드 10 · 증거금 표를 한 줄씩 재구성하기}
첫날 잔액 $1,892$에서 $400$ 손실이 나면 $1,492$로 유지증거금 $1,400$보다 높다. 다음날 $200$ 손실이 나면 $1,292$가 되어 유지증거금 아래로 내려간다. 초기 수준으로 복구해야 한다는 규칙이라면 납입액은 $1,892-1,292=\$600$이다. 이후 $2.13$에서 $1.93$으로 한 번에 내려가면 $20,000\times(-0.20)=-\$4,000$이므로 잔액은 일시적으로 음수가 될 수 있고, 초기 수준으로 되돌릴 더 큰 납입액이 생긴다.\par

\[
B_t^{\rm before\ call}=B_{t-1}^{\rm after\ call}+m q(F_t-F_{t-1}),\quad
C_t=\begin{cases}IM-B_t^{\rm before\ call},&B_t^{\rm before\ call}<MM,\\0,&\text{그 밖의 경우.}\end{cases}
\]

표의 날짜·가격은 정산 절차를 연습하기 위한 예시다. \textbf{누적 손익}과 \textbf{증거금 납입액}을 더해 투자수익으로 계산하면 안 된다. 납입액은 소유 계좌로 옮긴 현금이며, 계약의 손익 자체는 가격 변화에서 나온다.\par

초기 잔액 $1,892$, 유지 기준 $1,400$, 총수량 $20,000$부셸인 예를 날짜별로 쓰면 다음과 같다. 표에서 ‘정산 직후’는 추가 납입 전 잔액이며, 기준 미달 때 납입액을 계산한다.\par

\begin{center}\footnotesize\setlength{\tabcolsep}{3pt}\begin{longtable}{p{0.13\linewidth}p{0.12\linewidth}p{0.15\linewidth}p{0.19\linewidth}p{0.14\linewidth}p{0.18\linewidth}}\hline
시점 & 선물가격 & 그날 손익 & 정산 직후 잔액 & 추가 납입 & 납입 후 잔액 \\ \hline
계약 직후 & $2.21$ & — & $1,892$ & — & $1,892$ \\ 
1일 & $2.19$ & $-400$ & $1,492$ & $0$ & $1,492$ \\ 
2일 & $2.18$ & $-200$ & $1,292$ & $600$ & $1,892$ \\ 
3일 & $2.13$ & $-1,000$ & $892$ & $1,000$ & $1,892$ \\ 
4일 & $1.93$ & $-4,000$ & $-2,108$ & $4,000$ & $1,892$ \\ 
\hline\end{longtable}\end{center}\normalsize

가격이 $2.21$에서 $1.93$으로 변한 총 계약손실은 $20{,}000(1.93-2.21)=-\$5{,}600$이다. 추가 납입액의 합 $5{,}600$과 우연히 같지만, 전자는 \textbf{가격위험의 결과}, 후자는 그 손실 때문에 계좌 잔액을 유지하려고 옮긴 \textbf{현금의 합}이다. 담보를 차입해 조달했다면 이자비용은 별도로 발생한다.\par

\subsection{슬라이드 11 · 미결제약정·거래량·정산가격}
새 매수자와 새 매도자가 만나면 거래량은 1 늘고 미결제약정도 1 늘 수 있다. 기존 매수자와 기존 매도자가 서로 포지션을 청산하면 거래량은 1 늘지만 미결제약정은 1 줄 수 있다. 한쪽 신규·한쪽 기존이면 거래량만 늘고 미결제약정은 그대로일 수 있다. 정산가격은 일일손익과 담보를 계산하는 기준가격이며 그날 마지막 개별 체결가격과 반드시 같지는 않다.\par

\subsection{슬라이드 12 · 반대매매·실물인도·현금결제}
만기 전에 반대 방향 선물을 거래해 포지션을 닫으면 실물인도 의무를 대개 피할 수 있다. 실물인도가 가능한 계약에는 첫 인도통지일, 최종 인도통지일, 마지막 거래일이 따로 있어 \textbf{매매를 마칠 수 있는 날과 인도 절차가 시작되는 날}을 구별해야 한다. 현금결제형 지수선물은 만기의 최종 기준값과 선물 정산가격 차액으로 정리한다. 결제방식이 달라도 만기 직전 가격 수렴의 논리는 기초자산·계약 규격에 달려 있다.\par

\subsection{슬라이드 13 · 국내 선물상품을 읽는 공통표}
국내에는 주가지수, 개별주식, 금리, 통화, 원자재 등 여러 기초자산의 선물이 있다. 상품 이름 뒤에 \textbf{기초자산, 계약당 수량 또는 승수, 호가단위, 결제월, 최종결제방법}을 붙여 읽자. 예컨대 지수선물의 가격이 1포인트 움직여도 승수가 다르면 원화 손익이 다르다. 실제 거래 조건은 변경될 수 있으므로 계약을 적용할 때는 해당 거래소의 최신 명세를 확인한다.\par

\subsection{슬라이드 14 · KOSPI200 선물: 지수 포인트를 원화로 바꾸기}
KOSPI200 선물 한 계약의 명목가치는 \texttt{선물지수 × 승수}다. 승수가 25만 원이고 선물지수가 1,100이라면 2억 7,500만 원이다. 가격이 $0.05$포인트 움직이는 한 틱의 손익은 $0.05\times250,000=12,500$원이다. 선물가격이 오른 경우 매수자는 이익, 매도자는 손실이다. 승수가 5만 원인 미니 계약에서는 같은 지수 움직임의 원화 손익이 5분의 1이다. \href{https://global.krx.co.kr/contents/GLB/02/0201/0201040201/GLB0201040201.jsp}{한국거래소 상품명세}.\par

승수의 단위는 \textbf{원/지수포인트}다. 그래서 $1{,}100$포인트에 $250{,}000$원/포인트를 곱하면 원화 명목금액이 되고, $0.05$포인트에 같은 승수를 곱하면 한 틱의 원화 손익이 된다. 10계약을 매수한 뒤 지수가 1포인트 하락하면 손실은 $10\times250{,}000\times(-1)=-2{,}500{,}000$원이다. 명목금액 $27.5$억 원을 모두 지급한 것이 아니라 작은 증거금으로 이 가격변화에 노출된 것이므로, 손익률을 말할 때는 \textbf{명목금액 대비인지, 예치 담보 대비인지} 분모를 밝혀야 한다.\par

\par\medskip\noindent\textbf{심화: 승수와 증거금률을 함께 읽기}\quad 명목가치 2억 7,500만 원에 초기증거금률을 21.75\%로 가정하면 필요한 담보는 약 5,981만 원이다. 단순 명목/증거금 비율은 약 4.60배다. 실제 증거금률과 상품 규격은 시점과 계좌에 따라 달라지므로 거래소·중개회사의 적용 기준을 확인해야 한다.\par

\subsection{슬라이드 15 · 국채선물: 가격보다 금리 민감도를 먼저 보기}
국채선물의 기초는 표준화된 국채 또는 인도 가능한 국채 바스켓의 가격이다. 국채 가격과 금리는 반대로 움직이므로 \textbf{금리 상승을 헤지하려는 채권 보유자}는 보통 국채선물을 매도한다. 두 채권의 명목금액이 같아도 듀레이션이 다르면 금리 변화 손익은 다르다. 헤지 계약 수를 정할 때는 $\text{채권 DV01}/\text{선물 DV01}$처럼 1bp 금리변화당 가치변화를 맞추는 접근이 필요하다. 실제 결제방식과 계약 명세는 상품·시점에 따라 확인해야 한다.\par

$D_P$를 채권 포트폴리오 가치가 금리 1bp 상승에 따라 감소하는 원화 금액, $D_F$를 선물 1계약을 \textbf{매수}했을 때 같은 상승으로 잃는 원화 금액이라고 하자. 금리가 $\Delta y$bp 오르면 채권 손익은 대략 $-D_P\Delta y$이고, 선물 $N$계약을 \textbf{매도}한 손익은 $+ND_F\Delta y$다. 두 1차 효과를 상쇄하려면\par

\[
-D_P\Delta y+ND_F\Delta y=0
\quad\Longrightarrow\quad N=\frac{D_P}{D_F}.
\]

예를 들어 $D_P=1{,}000$만 원/bp, $D_F=50$만 원/bp라면 선물 20계약 매도가 출발점이다. 이는 금리곡선 전체가 같은 폭으로 움직인다는 \textbf{평행이동 근사}다. 포트폴리오와 인도채의 만기가 다르거나 최저가인도채권이 바뀌면 잔여 위험이 생기므로 실제 헤지는 전환계수와 금리구간별 민감도를 확인한다.\par

\subsection{슬라이드 16 · 달러선물: 환율 표시 방향과 계약 크기}
원/달러 환율은 1달러의 원화가격이다. 달러선물 매수자는 환율 상승에 이익을 얻으므로 미래에 달러를 사야 하는 수입업자의 원화 비용 증가를 상쇄할 수 있다. 예시에서 계약당 달러 수량이 1만 달러이고 가격이 1원 오르면 한 계약의 손익은 1만 원이다. 가격 표기, 호가단위, 최종결제방법을 확인하지 않고 미국달러선물과 다른 환율 선물을 같은 것으로 다루면 단위 오류가 생긴다.\par

\subsection{슬라이드 17 · 증거금 예시의 계산 검산}
선물지수 $232.5$와 승수 25만 원이면 계약 명목가치는 $232.5\times250,000=58,125,000$원이다. 초기증거금률 13.8\%라면 $8,021,250$원, 유지증거금률 9.2\%라면 $5,347,500$원이다. 이 값은 계약가격을 확정하는 ‘선도 가격’이 아니라 \textbf{담보 계산의 기준이 되는 선물 명목가치}다. 비율이 조정되면 요구액도 달라진다.\par

\subsection{슬라이드 18 · 미국 주가지수선물과 미니 계약}
미국 지수선물도 지수 포인트에 계약 승수를 곱해 원화 대신 달러 손익을 계산한다. E-mini처럼 계약 크기를 작게 설계하면 같은 지수 방향을 더 작은 명목노출로 거래할 수 있다. 하지만 작은 계약 10개를 사면 큰 계약과 같은 총노출이 생길 수 있으므로 상품명에 ‘미니’가 들어간다는 이유만으로 위험이 작다고 보면 안 된다. 거래소의 거래시간·최종정산 방법도 함께 확인해야 한다.\par

\subsection{슬라이드 19 · 미국 국채선물의 인도 선택권}
국채선물은 인도 가능한 여러 채권 중 실제로 어떤 채권을 넘길지 정하는 규칙이 있을 수 있다. 채권마다 쿠폰·만기가 다르므로 전환계수와 최저가인도채권(CTD)의 개념이 중요해진다. 지수선물처럼 단일 숫자에 승수만 곱하는 상품보다 기초자산의 선택과 금리곡선 움직임이 헤지 효과를 복잡하게 만든다. 입문 단계에서는 먼저 ‘채권 가격 상승 ↔ 금리 하락’의 방향을 확실히 익힌다.\par

\subsection{슬라이드 20 · 연방기금금리선물은 금리를 가격으로 표시한다}
금리선물은 가격을 \texttt{100 - 예상 평균 금리(\%)}처럼 표시하는 경우가 있다. 지표 금리 기대가 4\%에서 3.5\%로 내려가면 선물가격은 96에서 96.5로 오른다. 따라서 채권과 유사하게 \textbf{금리 하락에 선물가격 상승}이라는 방향이 나온다. 하지만 계약의 실제 지급액은 평균 금리 산정 기간·호가단위·승수에 달려 있으므로 계약 명세를 확인한다.\par

금리와 선물가격의 변화를 식으로 쓰면 $F=100-R$이므로 $\Delta F=-\Delta R$이다. 변동금리로 차입한 기업은 금리 상승 때 이자지급액이 증가한다. 이 기업이 금리선물을 매도했다면 $R$ 상승으로 $F$가 하락할 때 선물에서 이익을 얻어 증가한 이자비용을 일부 상쇄한다. 반대로 변동금리 수취자는 금리 하락을 걱정한다면 선물 매수가 자연스럽다. \textbf{금리 숫자의 방향과 호가된 선물가격의 방향이 반대}라는 사실을 먼저 적어야 포지션 부호가 헷갈리지 않는다.\par

\subsection{슬라이드 21 · 만기 인도보다 반대매매가 많은 이유}
선물로 가격위험만 이전하고 싶은 참가자는 실제 물품을 받을 필요가 없다. 만기 전 반대매매를 하면 가격 변화 손익은 일일정산으로 이미 실현되고 인도 절차는 피할 수 있다. 과거 상품별 실물인도·현금결제 비중이 낮았다는 통계는 \textbf{선물가격이 실물거래와 무관하다}는 뜻이 아니다. 인도가 가능한 계약에서는 소수의 인도 가능성이 만기 가격의 현물 수렴을 지탱한다.\par

\subsection{슬라이드 22 · 선도와 선물의 같은 점부터}
둘 다 미래 거래가격을 오늘 정해 기초자산 가격에 선형적으로 노출된다. 매수자는 기초가격 상승에 유리하고 매도자는 하락에 유리하다. 아래 차이는 지급액의 모양보다 \textbf{계약 표준화와 중간 정산 현금흐름}에서 생긴다. 따라서 ‘선물은 선도의 다른 이름’으로 외우면 신용·유동성 문제를 놓친다.\par

\subsection{슬라이드 23 · 선도는 맞춤형 양자 약속}
선도는 수량 $Q$, 인도가격 $K$, 만기 $T$, 결제통화와 실물/현금결제 조건을 두 당사자가 정한다. 만기에 매수의 차액결제 손익은 $Q(S_T-K)$다. 실물인도면 매수자는 $QK$를 내고 자산 $Q$개를 받으므로 동일한 경제적 순손익이 된다. 계약 종료 전의 가치는 시장조건에 따라 달라지고, 만기 지급 전에 거래상대방의 부도 위험이 존재한다.\par

\subsection{슬라이드 24 · 선물은 표준화와 일일정산}
선물도 $Q(S_T-K)$에 가까운 누적 가격변화 손익을 갖지만 일별 가격 변화가 매일 현금화된다. 초기증거금은 계약의 구입비가 아니라 이행 담보이고, 추가 증거금은 가격이 불리할 때 현금을 요구한다. 청산소가 양자 신용위험을 줄이는 대신 \textbf{현금흐름의 타이밍}이 선도와 달라진다. 이 차이가 금리와 선물가격 변화가 관련될 때 두 가격의 작은 차이를 만들 수 있다.\par

\subsection{슬라이드 25 · 비교표의 절대적인 표현을 점검하기}
장외 선도가 반드시 담보 없이 거래되는 것은 아니다. 신용지원약정(CSA)에 따라 담보를 주고받을 수 있다. 반대로 거래소 선물도 청산회원·청산소의 극단적 실패 가능성이 완전히 0은 아니다. 선도도 중도해지·양도·반대 계약으로 경제적 노출을 줄일 수 있고, 선물 중에도 만기 실물인도형이 있다. \textbf{‘보통 그렇다’는 제도적 경향과 ‘항상 그렇다’는 수학적 성질}을 구별하자.\par

\subsection{슬라이드 26 · 표로 정리하는 핵심 차이}
\begin{center}\small\begin{longtable}{p{0.22\linewidth}p{0.32\linewidth}p{0.32\linewidth}}\hline
항목 & 선도 & 선물 \\ \hline
조건 & 필요에 맞춰 약정 & 거래소 규격 \\ 
주된 상대방 구조 & 양자 약속, 필요시 담보 & 청산소와 회원 체계 \\ 
손익 실현 시점 & 보통 만기 또는 중도종료 & 매일 정산 \\ 
초기 현금 & 일반적인 신규계약 가치는 0, 담보는 별개 & 초기증거금 필요 \\ 
주된 추가 위험 & 거래상대방 신용위험, 유동성 & 증거금 현금부족, 베이시스 \\ 
\hline\end{longtable}\end{center}\normalsize

실제 계약에는 예외가 있으므로 표는 분석의 시작점이다. 어느 쪽이 언제 유리한지는 거래자의 기존 노출과 자금조달 능력에 달린다.\par

\subsection{슬라이드 27 · ‘선도가격’과 기존 계약가치의 구분}
오늘 새로 체결하는 공정 선도 인도가격을 $F_{0,T}$라 한다. 공정 가격으로 시작하면 보통 계약가치 $f_0=0$이다. 시간이 지나 현물가격이 바뀌면 \textbf{기존 계약의 인도가격 $K=F_{0,T}$는 그대로}, 새로운 공정 선도가격 $F_{t,T}$는 달라진다. 그래서 기존 매수 계약의 가치는 일반적으로 $f_t=e^{-r(T-t)}(F_{t,T}-K)$로 평가한다. 이 식은 다음 장에서 유도한다.\par

\subsection{슬라이드 28 · 옥수수 선도의 실물인도와 손익}
1백만 달러를 부셸당 $2.20$에 지급하기로 했다면 받을 양은 $Q=1,000,000/2.20\approx454,545.45$부셸이다. 만기 현물가격이 $2.10$이면 받은 옥수수의 시장가치는 약 $954,545.45$다. 약정 지급액 1백만 달러와의 차이인 약 $45,454.55$가 매수자의 손실이다. 현물가격이 내려가 현물시장에서 더 싸게 살 수 있었던 \textbf{기회비용}으로도 읽을 수 있다.\par

\[
\Pi_T^{\rm long}=Q(S_T-K)=\frac{1{,}000{,}000}{2.20}(2.10-2.20)\approx-45{,}454.55.
\]

\subsection{슬라이드 29 · 결제와 위험의 체크리스트}
포지션이 매수인지 매도인지, 계약 승수와 통화가 무엇인지, 정산가격이 언제 정해지는지, 유지증거금 아래로 내려갔을 때 얼마를 언제 납입해야 하는지를 순서대로 계산한다. 선도라면 거래상대방 신용위험과 중도종료가치도 적는다. 만기 손익이 같더라도 \textbf{결제 경로가 다른 두 계약의 실무 위험은 같지 않다.}\par

\subsection{슬라이드 30 · 다음 단계: 왜 그 선도가격인가}
이제 계약 구조를 알았으므로 오늘의 현물가격 $S_0$에서 만기 인도가격 $F_{0,T}$를 유도할 수 있다. 핵심은 현물을 지금 사서 자금조달·보관·배당을 처리한 거래와, 선도로 미래에 같은 자산을 받는 거래의 만기 현금흐름을 비교하는 것이다. 다음 장에서는 복리와 현재가치부터 시작한다.\par

\section{05. 선도가격, 보유비용, 베이시스}
가격을 예측하지 않아도 선도가격을 계산할 수 있다. 오늘 현물을 사서 만기까지 보유하는 비용과, 만기에 선도로 같은 자산을 받는 비용이 다르면 동일한 현금흐름의 가격 차이가 생긴다. 이 장은 복리에서 시작해 배당·쿠폰·외화금리·보관비용을 차례로 더한다.\par

\subsection{슬라이드 01 · 무차익 가격의 질문}
미래 주가가 오를지 내릴지와 오늘의 \textbf{공정 인도가격}은 다른 문제다. 투자자가 현물 매수와 선도 매수를 통해 만기에 같은 자산을 받을 수 있다면, 두 경로의 오늘 비용을 맞춰야 한다. 그 결과 얻는 $F_{0,T}$는 기대 현물가격 $E[S_T]$와 일반적으로 같지 않다. 위험 프리미엄과 투자자의 기댓값은 나중에 별도로 다룬다.\par

\subsection{슬라이드 02 · 이산복리에서 시작하기}
초기금액 $W_0$를 연 $r$의 명목금리로 $T$년간, 1년에 $m$번 복리로 운용한다고 하자. 한 번의 이자기간은 $1/m$년이고 매 기간 성장배수는 $1+r/m$다. 총 $mT$번 곱하므로\par

\[
W_T=W_0\left(1+\frac{r}{m}\right)^{mT}.
\]

$m=1$이면 연복리, $m=2$이면 반기복리다. 같은 명목금리 $r>0$라면 이자에 이자가 더 자주 붙을수록 만기금액이 커진다.\par

\subsection{슬라이드 03 · 복리 빈도를 숫자로 확인하기}
$W_0=1,000$, $r=10\%$, $T=5$년이면 연복리는 $1,000(1.1)^5\approx1,610.51$이다. 반기복리는 $1,000(1.05)^{10}\approx1,628.89$, 월복리는 $1,000(1+0.1/12)^{60}\approx1,645.31$이다. 이는 \textbf{같은 명목 연율}을 서로 다른 빈도로 복리한 결과이지, 서로 다른 위험을 비교한 것이 아니다. 현금흐름이 언제 지급되는지에 맞는 금리 관례를 써야 한다.\par

\subsection{슬라이드 04 · 연속복리와 할인}
$m$을 무한히 크게 할 때 $\lim_{m\to\infty}(1+r/m)^{mT}=e^{rT}$이므로 $W_T=W_0e^{rT}$다. 위 숫자는 $1,000e^{0.5}\approx1,648.72$가 된다. 확정된 미래금액 $X_T$를 오늘 얼마로 바꿀지는 이 식을 뒤집어 $PV_0(X_T)=X_Te^{-rT}$로 구한다. 시간 0의 1원과 시간 $T$의 1원을 같은 칸에 더하지 않는 것이 모든 선도가격 유도의 출발점이다.\par

\[
r_{\rm cc}=\frac{1}{T}\ln\frac{W_T}{W_0},\qquad
PV_0(X_T)=X_Te^{-r_{\rm cc}T}.
\]

\par\medskip\noindent\textbf{심화: 서로 다른 만기 금리를 쓸 때}\quad 6개월 현금흐름에는 6개월 무위험 할인율, 1년 현금흐름에는 1년 무위험 할인율을 쓴다. $r_{0,t}$를 만기 $t$의 연속복리 무위험 금리로 두면 할인인자는 $e^{-r_{0,t}t}$다. 모든 만기에 동일한 $r$을 쓰는 식은 설명을 위한 단순화다.\par

\subsection{슬라이드 05 · $S$, $F$, $K$, $f$를 분리하기}
$S_t$는 시점 $t$의 현물가격, $F_{t,T}$는 시점 $t$에 새로 약정할 만기 $T$의 공정 선도가격이다. $K$는 이미 체결한 계약에 적힌 \textbf{고정 인도가격}이며, $f_t$는 그 오래된 계약의 현재가치다. 새 계약은 공정 $K=F_{t,T}$로 시작해 보통 $f_t=0$이다. 시간이 흐른 뒤 $F_{t,T}$가 움직여도 기존 $K$는 바뀌지 않으므로 $f_t$가 0이 아니게 된다.\par

\subsection{슬라이드 06 · 투자자산과 소비자산}
주식·채권처럼 보유자가 주로 투자 수익을 얻기 위해 들고 있고 공매도·대출이 가능한 자산은 현물과 선도의 현금흐름을 비교적 쉽게 복제한다. 석유나 구리처럼 생산에 즉시 투입할 수 있는 자산은 재고 보유의 \textbf{편익수익률}이 있어 부족한 재고를 공매도하기 어려울 수 있다. 둘 다 선도가격에 자금조달·보관비용이 영향을 주지만, 소비자산에는 가격을 한 점으로 고정하기 어려운 제약이 있다.\par

\subsection{슬라이드 07 · 보유비용 공식의 한 단계씩 유도}
배당이 없고 보관비용도 없는 투자자산 한 단위를 오늘 $S_0$에 사기 위해 $S_0$를 빌린다. 만기에 자산을 선도 매도해 $F_{0,T}$를 받고 차입금 $S_0e^{rT}$를 갚는다. 아무 위험 없이 남는 금액이 생기면 가격 차익이므로 공정가격은 $F_{0,T}=S_0e^{rT}$다. 선도가격이 이보다 높으면 \textbf{현물 매수·선도 매도}, 낮으면 자산 공매도와 선도 매수의 역거래를 고려한다. 역거래 가능성은 공매도와 차입 조건에 달린다.\par

배당 $I_i$를 $t_i<T$에 받는다면 현물 보유의 순비용은 $S_0-\sum_i I_i e^{-r_i t_i}$다. 미리 알 수 있는 보관비용 $C_j$가 있으면 그 현재가치를 더한다. 따라서 공정 선도가격은\par

\[
F_{0,T}=\left(S_0+\sum_j PV_0(C_j)-\sum_i PV_0(I_i)\right)e^{r_T T}.
\]

여기서 앞 괄호는 \textbf{오늘 자산을 마련해 만기까지 유지하는 순비용}이다. 이산 현금흐름과 비례적 연속 보관비용을 동시에 쓸 때는 각 비용을 중복계산하지 않도록 어느 항목에 포함했는지 명시한다.\par

\subsection{슬라이드 08 · 기존 선도계약의 현재가치}
시점 $t$에 남은 기간을 $\tau=T-t$라 하자. 기존 매수 계약과 오늘 체결하는 새 매도 계약을 같은 만기로 결합하면 만기에는 자산 인수·인도가 상쇄되어 \textbf{확정금액} $F_{t,T}-K$만 남는다. 그러므로\par

\[
f_t^{\rm long}=(F_{t,T}-K)e^{-r_t\tau}.
\]

배당이 없는 주식이라면 $F_{t,T}=S_te^{r_t\tau}$를 대입해 $f_t=S_t-Ke^{-r_t\tau}$가 된다. 알려진 배당의 현재가치가 $I_t$라면 $f_t=S_t-I_t-Ke^{-r_t\tau}$. 선도가격 $F_{t,T}$는 \textbf{새 약정의 인도가격}, $f_t$는 \textbf{기존 약정의 지금 청산 가치}라는 구별이 이 식에서 선명해진다.\par

\subsection{슬라이드 09 · 배당 없는 주식: 과대평가된 선도}
90일을 $T=1/4$년으로, $S_0=\$20$, 연속복리 금리 $r=4\%$로 둔다. 무차익 선도가격은 $20e^{0.04/4}\approx\$20.2010$이다. 시장 선도가격이 $21$이면 너무 비싸다. 오늘 $20$를 빌려 주식을 사고, 동시에 만기에 그 주식을 $21$에 팔기로 약정한다. 오늘 순현금흐름은 0이다.\par

\subsection{슬라이드 10 · 만기 현금흐름이 왜 확정되는가}
만기에 주식의 시장가치는 $S_T$지만 선도 매도에 주식을 넘기고 $21$을 받으므로 $S_T$는 사라진다. 빚 상환액은 $20e^{0.01}\approx20.2010$이다. 확정 이익은 $21-20.2010\approx\$0.7990$다.\par

\begin{center}\small\begin{longtable}{p{0.22\linewidth}p{0.32\linewidth}p{0.32\linewidth}}\hline
거래 & 오늘 & 만기 \\ \hline
$20$ 차입 & $+20$ & $-20e^{0.01}$ \\ 
주식 1주 매수 & $-20$ & $+S_T$ \\ 
선도 1주 매도 & $0$ & $21-S_T$ \\ 
합계 & $0$ & $21-20e^{0.01}$ \\ 
\hline\end{longtable}\end{center}\normalsize

거래비용이나 공매도·자금조달 제약이 있는 현실에서는 정확히 한 점이 아니라 차익거래가 작동하지 않는 \textbf{가격 범위}가 생길 수 있다.\par

\subsection{슬라이드 11 · 이표채 선도: 쿠폰 시점에 맞춰 할인하기}
현재 채권가격 $1,200$, 6개월과 1년의 쿠폰이 각각 $50$, 인도 직후에는 1년 쿠폰을 선도 매수자가 받지 않는 \textbf{배당락·쿠폰락 기준}이라고 하자. 6개월 금리 8\%, 1년 금리 10\%를 연속복리로 쓰면 선도 매수자가 못 받는 쿠폰의 현재가치는 $50e^{-0.08/2}+50e^{-0.10}$이다. 따라서\par

\[
F_{0,1}=\bigl(1{,}200-50e^{-0.04}-50e^{-0.10}\bigr)e^{0.10}
\approx1{,}223.11.
\]

만기 쿠폰의 귀속이 다르면 마지막 $50$의 처리도 달라진다. \textbf{인도 시점과 쿠폰 지급 시점의 선후}를 계약에서 확인해야 한다.\par

\subsection{슬라이드 12 · 이표채 차익거래를 현금흐름으로 검산}
현물채권을 사기 위해 $48.04$를 6개월까지 빌리고 나머지 $1,151.96$을 1년까지 빌린다. 6개월에는 첫 쿠폰 $50$으로 첫 빚을 갚는다. 1년에는 채권의 두 번째 쿠폰 $50$과 선도 매도 인도대금 $F$를 받고, 장기 차입금 $1,151.96e^{0.10}\approx1,273.11$을 갚는다. 순이익은 $F+50-1,273.11=F-1,223.11$이다. 시장가격 $F=1,230$이면 약 $6.89$의 확정 이익이 생긴다. 이 계산이 앞 슬라이드의 식을 독립적으로 검증한다.\par

\subsection{슬라이드 13 · 외화는 ‘보유하면 외국 금리를 받는 자산’이다}
환율 $S_0$를 달러/마르크(DM), 즉 \textbf{1 DM의 달러 가격}으로 표시한다. 달러 금리 $r_{\$}=5\%$, 마르크 금리 $r_{DM}=8\%$라면 마르크를 보유해 외국 금리를 받을 수 있다. 따라서 외국 금리가 배당수익률처럼 달러 기준 선도가격을 낮춘다. 60일을 예시처럼 $T=1/6$년으로 놓으면\par

\[
F_{0,T}=S_0e^{(r_{\$}-r_{DM})T}
=0.55e^{(0.05-0.08)/6}\approx0.54726\ \text{달러/DM}.
\]

역사적 통화 DM은 유로 도입 전의 예시다. 오늘의 실제 환율이나 금리를 나타내는 숫자로 읽지 않는다.\par

\subsection{슬라이드 14 · 커버드 금리평가의 복제 거래}
만기에 1 DM을 확보하려면 오늘 $e^{-r_{DM}T}$ DM을 사서 독일에서 운용한다. 그 비용은 달러로 $S_0e^{-r_{DM}T}$이고, 달러 차입으로 조달하면 만기 상환액은 $S_0e^{(r_{\$}-r_{DM})T}$다. 다른 방법은 오늘 1 DM 매수 선도를 체결해 만기에 $F_{0,T}$달러를 지불하는 것이다. 둘 다 만기에 정확히 1 DM을 주므로 비용이 같아야 한다. \href{https://www.bis.org/publications/qr-201609/covered-interest-parity-lost-understanding-cross-currency-basis}{BIS의 커버드 금리평가 설명}.\par

현실에서는 달러 조달비용·규제 자본·상대방 위험 때문에 이 단순 등식과 시장가격 사이에 \texttt{cross-currency basis}가 지속될 수 있다. 이는 식이 쓸모없다는 뜻이 아니라, \textbf{복제 거래에 필요한 조건}이 비용 없이 충족되지 않는다는 신호다.\par

\subsection{슬라이드 15 · 주가지수의 연속 배당수익률}
기초지수가 연속 배당수익률 $q$를 지급하고 배당을 재투자할 수 있다고 하자. 만기에 1단위의 지수를 확보하려면 오늘 $e^{-qT}$단위를 사면 되고, 그 비용은 $S_0e^{-qT}$다. 이를 금리 $r$로 조달하면 만기비용은 $S_0e^{(r-q)T}$. 따라서\par

\[
F_{0,T}=S_0e^{(r-q)T}.
\]

배당수익률이 시간에 따라 알려진 함수 $q(u)$라면 $qT$ 대신 $\int_0^Tq(u)\,du$를 쓴다. 알려진 개별 배당금에는 슬라이드 11의 \textbf{이산 배당 현재가치 식}이 더 직접적이다. \href{https://www.cmegroup.com/trading/equity-index/fairvalue.html}{CME 지수선물 공정가치 설명}.\par

\subsection{슬라이드 16 · 소비자산과 편익수익률의 범위}
원유를 보관하면 자금조달비용과 창고비용이 들지만, 재고가 있으면 생산 중단을 피하거나 일시적 부족을 이용할 수 있다. 이 \textbf{재고의 편익}을 관측하기 어려운 편익수익률 $y$로 나타내면 설명식으로 $F\approx S_0e^{(r+s-y)T}$를 쓴다. 그러나 이것은 편익을 모두 확정 현금흐름처럼 복제한 엄밀한 차익거래 등식이 아니다.\par

현물을 살 수 있고 보관할 수 있다면 선도 과대평가에는 현물 매수·선도 매도로 대응하므로 $F\leq S_0e^{(r+s)T}$라는 상한이 생길 수 있다. 반대 방향에는 현물 공매도가 어렵고 재고를 포기할 수 없어 같은 강도의 하한이 없을 수 있다. \textbf{투자자산의 등식과 소비자산의 경계 조건}을 구분해야 한다.\par

\subsection{슬라이드 17 · 과대·과소평가를 발견했을 때의 거래}
시장 선도가격이 복제비용보다 높으면 선도를 매도하고 현물을 사서 보유한다. 낮으면 선도를 매수하고 현물을 공매도한 뒤 공매도 대금을 운용한다. 후자의 거래는 대주 가능 여부, 배당 지급 의무, 재고 보유자의 편익 때문에 어려울 수 있다. 선물은 일일정산이 있어 여기에 증거금의 현금흐름과 차입비용도 더해진다. ‘이론값과 다르다’와 ‘오늘 누구나 확정 이익을 낸다’는 같은 명제가 아니다.\par

\subsection{슬라이드 18 · 지수선물의 이론값 계산}
현물지수 $S_0=8,300$, 연속복리 금리 $r=8\%$, 배당수익률 $q=5\%$, 만기 $T=0.5$년이면\par

\[
F^*_{0,T}=8{,}300e^{(0.08-0.05)\times0.5}
=8{,}300e^{0.015}\approx8{,}425.44.
\]

시장 선도가격 $8,494$는 약 $68.56$포인트 높다. 지수 1단위의 현금복제가 가능하고 거래비용이 없다면 과대평가다. 계산의 관건은 금리 8\%에서 배당 5\%를 \textbf{빼는 이유}다. 지수를 보유하면 배당을 받기 때문에 미래 인도를 위해 오늘 지수를 사는 순비용이 줄어든다.\par

\subsection{슬라이드 19 · 지수 차익거래의 단위별 현금흐름}
오늘 $e^{-0.05\times0.5}\approx0.97531$단위의 지수를 사면 배당 재투자로 만기에 1단위가 된다. 비용은 약 $8,095.07$이며 이를 차입하면 만기 상환액은 $8,095.07e^{0.08\times0.5}\approx8,425.44$다. 선도 매도로 만기에 지수 1단위를 $8,494$에 넘기면 약 $68.56$의 확정 차액이 남는다. 실제 거래에는 지수 바스켓 구성·배당 예측·거래비용·차입 조건이 포함되어 오차가 생긴다.\par

\subsection{슬라이드 20 · 서로 다른 두 ‘콘탱고/백워데이션’}
첫째, \textbf{현물과 선도 비교}에서는 $F_{0,T}>S_0$를 콘탱고, $F_{0,T}<S_0$를 백워데이션이라 부른다. 둘째, \textbf{미래 현물 기대값과 비교}하는 ‘정상 백워데이션’은 $F_{0,T}<E^{P}[S_T]$처럼 실제 확률 $P$ 아래의 기대와 선도가격을 비교한다. 두 기준은 다르므로 $F>S_0$이면서 동시에 $F<E^{P}[S_T]$일 수 있다.\par

이 가이드의 베이시스는 $b_t=F_{t,T}-S_t$로 정의한다. 다른 시장에서는 $S_t-F_{t,T}$를 베이시스로 부르기도 한다. 부호가 달라진다는 사실을 먼저 적어야 \texttt{베이시스 확대}라는 말을 정확히 해석할 수 있다. $F=S e^{(r-q)T}$라면 $b=S(e^{(r-q)T}-1)$이고, 만기가 가까워질수록 이상적 조건에서 0으로 수렴한다. \href{https://www.cmegroup.com/education/courses/introduction-to-grains-and-oilseeds/learn-about-grain-convergence}{CME의 현물·선물 수렴 설명}.\par

\subsection{슬라이드 21 · KOSPI200 그래프에서 읽을 수 있는 것과 없는 것}
현물지수와 선물가격의 차이는 만기, 금리, 예상 배당, 거래 수급, 제약에 따라 움직인다. 역사적 그래프에서 베이시스가 음수였다는 사실만으로 시장이 앞으로 떨어질 것이라고 결론 낼 수 없다. 현물과 선물의 표시 시점이 같은지, 같은 만기의 선물을 연결한 자료인지, 배당락일과 만기 교체가 반영됐는지를 확인해야 한다. 선물이 현물보다 먼저 움직이는 경우가 있어도 \textbf{선물이 현물을 항상 일방적으로 결정한다}는 뜻은 아니다.\par

\subsection{슬라이드 22 · 선물과 선도가격은 언제 달라지는가}
선물은 유리한 가격변화의 이익을 매일 현금으로 받고 불리한 손실을 매일 낸다. 선도는 보통 만기까지 한 번에 정산한다. 금리가 확정적이면 중간 정산액의 재투자·조달 조건이 예측 가능해 두 가격이 가까워진다. 금리가 확률적으로 변하고 선물가격 변화와 양의 상관관계가 있으면 이익이 나올 때 금리도 높아 재투자에 유리해 선물가격이 선도가격보다 높아질 수 있다. 반대 상관관계면 방향이 뒤집힐 수 있다. 이것은 \textbf{특정 모형·담보·재투자 가정 아래의 경향}이지 모든 시장에서 부호를 보장하는 법칙은 아니다.\par

\par\medskip\noindent\textbf{심화: 왜 중간 현금흐름이 가격에 영향을 주나}\quad 만기 손익이 같더라도 A계약은 오늘 1원을 받고 B계약은 만기 1원을 받는다면 두 계약의 현재가치는 다르다. 선물의 일일정산액은 그날 이후의 금리로 운용하거나 조달해야 한다. 이 미래 금리가 선물 손익과 함께 변하면 ‘손익의 합계’와 ‘그 손익을 운용한 만기 재산’이 달라진다.\par

\subsection{슬라이드 23 · 복제 가능성이 결론을 결정한다}
투자자산은 동일한 미래 현금흐름을 현물 보유와 차입·대여로 만들 수 있어 무차익 \textbf{등식}이 강하게 작동한다. 소비자산에는 재고 편익·공매도 제약으로 한쪽 차익거래만 가능해 \textbf{상한이나 가격 구간}을 얻는 경우가 많다. 현물-선물 베이시스의 부호와 선도-미래현물 기대값의 비교를 섞지 않으면 가격이 왜 그렇게 형성되는지 명확해진다.\par

\subsection{슬라이드 24 · 다음 장의 질문: 적정가격으로 위험을 얼마나 줄일까}
공정가격 $F_{0,T}$를 알게 되면 기존 현금흐름에 $Q(S_T-F_{0,T})$를 더할 수 있다. $Q$의 부호와 크기를 조절해 위험을 완전히 없애거나 일부만 남긴다. 이후에는 평균 현금흐름만이 아니라 \textbf{하위 5\%의 현금흐름}을 계산하고, 비선형 사업수익에서는 왜 선형 선도가 완전한 보험이 되지 못하는지 본다.\par

\section{06. 선도 헤지와 현금흐름의 하방위험}
공정한 선도가격을 안 뒤에도 ‘얼마나 헤지할 것인가’라는 의사결정은 남는다. 완전헤지는 환율 민감도를 0으로 만들지만 기대 현금흐름을 바꾸고, 부분헤지는 남은 위험과 기대값을 함께 조절한다. 사업수익이 환율에 따라 불연속적으로 변하면 선형 선도만으로는 전체 분포를 한 점에 모을 수 없다.\par

\subsection{슬라이드 01 · 가격에서 위험관리로}
앞 장의 무차익 가격 $F_{0,T}$는 \textbf{계약을 맺는 공정한 환율}을 알려 준다. 이 장의 변수 $Q$는 계약의 규모와 방향을 뜻한다. 어떤 기업의 기존 현금흐름을 $C(S_T)$라 하면 선도 매수 $Q$를 더한 최종 현금흐름은\par

\[
CF_Q(S_T)=C(S_T)+Q(S_T-F_{0,T}).
\]

$Q>0$은 기초통화 매수, $Q<0$은 매도다. 최종 위험은 $C$와 선도손익의 \textbf{합}으로 계산한다.\par

\subsection{슬라이드 02 · 헤지의 성과 지표를 먼저 정하기}
분산을 최소화할지, 5\% 하위 현금흐름을 목표 이상으로 올릴지, 특정 확률로 연구개발비를 확보할지에 따라 최적 $Q$가 달라진다. 예를 들어 $\Pr(CF_Q\geq220)\geq95\%$라는 목표는 ‘평균 현금흐름이 220 이상’보다 강하다. 평균은 좋은 상태의 큰 수익으로 올라갈 수 있어도 나쁜 상태의 자금 부족을 감출 수 있기 때문이다.\par

\subsection{슬라이드 03 · 헤지비율과 방향}
기존 사업이 환율 $S_T$ 상승에 이익을 본다면 선도 매도($Q<0$)가 그 이익을 줄이는 동시에 하락 손실을 상쇄한다. 선형 노출이 $A S_T$이면 $Q=-A$가 완전헤지다. 헤지비율 $h=-Q/A$를 정의하면 $h=0$은 미헤지, $h=1$은 완전헤지다. $0<h<1$은 부분헤지이며, $h>1$은 기존 노출보다 많이 매도해 위험 방향을 뒤집는 \textbf{과잉헤지}다. 모든 상품에서 ‘100\%’가 자동으로 최적이라는 뜻은 아니다.\par

\subsection{슬라이드 04 · 원/달러 선도 매수·매도 손익선}
달러 1달러의 약정 인도가격을 $K=1,026$원/달러라 하면 달러 선도 매수 손익은 $S_T-1,026$원, 매도 손익은 $1,026-S_T$원이다. 두 직선은 $S_T=1,026$에서 0과 만나고 기울기는 각각 $+1$, $-1$이다. 가로축의 단위가 원/달러이므로 ‘달러가 강해진다’는 것은 이 그림에서 오른쪽으로 움직인다는 뜻이다.\par

\subsection{슬라이드 05 · 수입업자의 구매비용과 부분헤지}
만기에 1달러가 필요한 기업의 미헤지 원화 비용은 $S_T$다. 달러 선도 매수로 $h$달러를 미리 확보하면 총 원화 비용은\par

\[
C_h(S_T)=S_T-h(S_T-K)=(1-h)S_T+hK.
\]

$h=0.5$면 기울기가 $0.5$, $h=0.75$면 $0.25$, $h=1$이면 비용이 $K$로 완전히 고정된다. 높은 환율 구간의 나쁜 결과를 줄이는 대신 낮은 환율 구간의 유리한 결과도 포기한다. 이를 선도 ‘프리미엄’이라고 부르면 오해가 생긴다. 선도 신규계약에 직접 내는 옵션 프리미엄은 보통 없지만 \textbf{기회비용}은 있다.\par

\subsection{슬라이드 06 · 분포가 왜 좁아지는가}
$S_T$의 표준편차를 $\sigma_S$라 하고 $K,h$가 확정이라면 $C_h=(1-h)S_T+hK$이므로 $\operatorname{Var}(C_h)=(1-h)^2\operatorname{Var}(S_T)$이다. 완전헤지에서 환율만이 위험이라면 분산은 0이다. 부분헤지에서는 표준편차가 $|1-h|\sigma_S$로 줄어든다. 그림의 종 모양이 좁아지는 이유를 이 식이 설명한다. 실제 비용에 물량·수요·베이시스 위험이 있으면 완전헤지 뒤에도 분산이 남는다.\par

\subsection{슬라이드 07 · 외화선도 한 건의 두 현금흐름}
기업이 1개월 뒤 1,000만 파운드를 1파운드당 1.5달러에 사기로 하면 계약 규모는 \textbf{파운드 1,000만}, 달러 지급액은 \textbf{1,500만 달러}다. 만기에는 파운드를 받고 달러를 낸다. 차액결제라면 달러 기준 손익은 $10{,}000{,}000(S_T-1.5)$이다. 실제 인도와 차액결제는 계약 조건에 따라 다르지만 환율 방향의 경제적 노출은 같다. \texttt{T+2} 같은 실제 자금 결제일은 계약 만기·관측일과 구별한다.\par

\subsection{슬라이드 08 · 환율의 분자와 분모를 바꾸면 포지션도 바뀐다}
$F=1.5$ USD/GBP는 \textbf{1파운드의 달러 가격}이다. 이 가격으로 파운드를 매수한다는 것은 GBP 매수·USD 매도다. 역환율 GBP/USD로 바꾸면 수치가 $1/1.5$가 되지만 선도 가격을 단순히 역수로 바꾸는 데는 매수·매도 호가와 결제 통화의 차이도 고려해야 한다. 계약 수량 $Q$가 어떤 통화 단위인지 적지 않으면 $Q(S_T-F)$의 단위가 모호해진다.\par

\subsection{슬라이드 09 · AnGen의 두 시장: 환율이 매출을 두 번 움직인다}
이 기업의 미국 매출은 기본 3억 달러이고, 일본 경쟁사의 달러 가격 변화에 대응해 환율에 따라 변한다. 일본 매출은 520억 엔으로, 달러 환산액은 환율 $S_T$(USD/JPY)에 비례한다. 금액을 \textbf{백만 달러}로 표현하면 520억 엔은 52,000백만 엔이므로 일본 매출은 $52{,}000S_T$백만 달러다. 미국 매출의 환율 민감도 40,000을 더하면 사업 전체 민감도는 92,000이다. 미국 매출의 기준값을 기대 환율에서 300으로 맞추어 두면 아래 식이 나온다.\par

\[
R_{US}=300+40{,}000(S_T-E[S_T]),\quad
R_{JP}=52{,}000S_T,\quad C=450.
\]

\subsection{슬라이드 10 · 로그정규 환율의 평균과 하위 5\%}
환율의 로그변화 $X=\ln(S_T/S_0)$가 평균 $\mu=0.011153$, 표준편차 $\sigma=0.05$인 정규분포라고 하자. 그러면 $S_T=S_0e^X$이고 $e^X$의 평균 성질 때문에\par

\[
E[S_T]=S_0e^{\mu+\sigma^2/2}.
\]

$S_0=0.007692$이면 $E[S_T]\approx0.007788$ USD/JPY다. 정규분포 하위 5\%의 표준점수 $z_{0.05}\approx-1.64485$를 쓰면\par

\[
s_{0.05}=S_0e^{\mu+\sigma z_{0.05}}
=0.007692e^{0.011153-1.64485(0.05)}\approx0.007164.
\]

여기서 $\mu$는 \textbf{로그환율 변화의 평균}, $E[S_T]$는 \textbf{환율 수준의 평균}이다. 둘을 같은 숫자로 두면 평균 현금흐름부터 틀린다.\par

\subsection{슬라이드 11 · 현금흐름 모형과 단위를 한 번 더 검산하기}
비연구개발 비용은 4억 5천만 달러이므로 단위를 백만 달러로 맞추면 $450$이다. 따라서 선도를 포함하지 않은 사업 현금흐름은\par

\[
CF_0(S_T)=300+40{,}000(S_T-E[S_T])+52{,}000S_T-450
=-150+92{,}000S_T-40{,}000E[S_T].
\]

$S_T$의 단위가 달러/엔, 계수의 단위가 백만 엔이므로 곱은 백만 달러다. 비용 $450$을 원화나 엔화로 착각하면 숫자가 수백 배 달라진다. 현금흐름 위험의 출발점은 $\partial CF_0/\partial S_T=92{,}000$백만 엔이다.\par

\subsection{슬라이드 12 · 미헤지와 완전헤지의 기대값·하위값}
현금흐름이 $S_T$에 대해 증가하므로 미헤지의 하위 5\% 현금흐름은 $s_{0.05}$에서 나온다. 계산하면\par

\[
E[CF_0]=52{,}000E[S_T]-150\approx254.976,\qquad
q_{0.05}(CF_0)=92{,}000s_{0.05}-40{,}000E[S_T]-150\approx197.583.
\]

완전헤지는 $Q=-92{,}000$백만 엔을 \textbf{매도}한다. $CF_Q=CF_0+Q(S_T-F)$에 넣으면 $S_T$ 항이 없어져\par

\[
CF_{-92{,}000}=92{,}000F-40{,}000E[S_T]-150\approx248.720
\]

백만 달러가 모든 상태에서 확정된다($F=0.00772$). 평균은 낮지만 5\% 하방 현금흐름은 높다. $E[S_T]>F$라는 숫자에서는 선도 매도로 평균 상승의 일부를 포기했기 때문이다. \textbf{공정 선도가격이 미래 현물의 실제 기대값과 같아야 한다는 법칙은 없다.}\par

\subsection{슬라이드 13 · 2억 2천만 달러를 확보하는 부분헤지}
$Q$를 백만 엔 단위의 선도 매수량으로 정의하면 매도는 음수다. $-92{,}000\leq Q\leq0$ 범위에서는 $CF_Q$가 환율에 따라 증가하므로 5\% 현금흐름을 $s_{0.05}$에 대입할 수 있다.\par

\[
q_{0.05}(CF_Q)=92{,}000s_{0.05}-40{,}000E[S_T]-150+Q(s_{0.05}-F).
\]

이를 목표 $220$에 맞춰 $Q=[220-q_{0.05}(CF_0)]/(s_{0.05}-F)\approx-40{,}330$을 얻는다. 즉 약 403.3억 엔의 매도 선도다. 기대 현금흐름은\par

\[
E[CF_Q]=52{,}000E[S_T]-150+Q(E[S_T]-F)\approx252.234
\]

백만 달러다. 분모 $s_{0.05}-F<0$이므로 하방 현금흐름을 올리려면 $Q<0$이어야 한다는 부호도 확인할 수 있다.\par

\subsection{슬라이드 14 · 평균과 하방의 교환관계}
미헤지 $Q=0$은 $(q_{0.05},E[CF])\approx(197.58,254.98)$, 부분헤지 $Q=-40{,}330$은 $(220,252.23)$, 완전헤지 $Q=-92{,}000$은 $(248.72,248.72)$이다. 가로축을 하위 5\% 현금흐름, 세로축을 평균으로 두면 헤지가 커질수록 왼쪽 아래에서 오른쪽 아래로 움직인다. 완전헤지 수량은 $-92{,}000$백만 엔, 즉 엔화 920억 엔의 매도이며, \textbf{수량과 통화 단위를 함께 적어야} 크기를 정확히 해석할 수 있다.\par

이 직선적인 교환관계는 현금흐름이 $S_T$의 아핀함수이고 $Q$가 $-92{,}000$ 이상일 때 성립한다. 과잉헤지하면 하위 5\%를 만드는 환율 상태가 \textbf{상위 5\% 환율}로 바뀌므로 같은 식을 계속 쓰면 안 된다.\par

\subsection{슬라이드 15 · 비선형 매출은 왜 선도로 고정되지 않는가}
이제 일본 매출 민감도가 환율 구간에 따라 달라진다고 하자. 미국 매출은 300백만 달러, 기존 비용은 450백만 달러다. 일본 매출 계수를 $a(s)$로 쓰면\par

\[
a(s)=\begin{cases}
10{,}000,&s\leq k_2=0.0065,\\
45{,}000,&k_2<s\leq k_1=0.0075,\\
55{,}000,&s>k_1,
\end{cases}\qquad CF_0(s)=-150+a(s)s.
\]

이 단순화한 수요모형은 임계값에서 불연속이다. $k_2$ 바로 아래와 위의 현금흐름은 각각 약 $-85$와 $142.5$백만 달러, $k_1$ 바로 아래와 위는 $187.5$와 $262.5$백만 달러다. 이 \textbf{수요의 점프}는 기울기가 일정한 선도 손익으로 상쇄할 수 없다. 연구개발 목표도 여기서는 200백만 달러로 바뀐다.\par

\subsection{슬라이드 16 · 비선형 현금흐름의 분포}
동일한 로그정규 환율을 사용하면 높은 환율 구간 $S_T>k_1$의 확률은 약 $76.688\%$, 중간 구간 $k_2<S_T\leq k_1$은 $23.295\%$, 낮은 구간 $S_T\leq k_2$는 $0.0165\%$다. 사업 현금흐름의 분포는 각 구간의 다른 직선과 임계값의 점프를 합친 것이므로 종 모양 하나로 근사하기 어렵다. 미헤지 평균은 약 $261.317$백만 달러이지만 하위 5\% 현금흐름은 약 $172.387$백만 달러다. 평균만 보면 목표 200을 넉넉히 넘지만 나쁜 상태의 자금은 부족하다.\par

\subsection{슬라이드 17 · 비선형 평균을 적분으로 구하기}
각 구간의 계수와 그 구간에서의 $S_T$ 평균 기여분을 곱해야 한다. $M(a)=E[S_T\mathbf1_{\{S_T\leq a\}}]$라 두면 로그정규분포의 적분을 완전제곱으로 정리해\par

\[
M(a)=S_0e^{\mu+\sigma^2/2}
\Phi\!\left(\frac{\ln(a/S_0)-\mu-\sigma^2}{\sigma}\right).
\]

이를 사용하면 선도 수량 $Q$를 포함한 평균은\par

\[
E[CF_Q]=-150+10{,}000M(k_2)+45{,}000[M(k_1)-M(k_2)]
+55{,}000[E[S_T]-M(k_1)]+Q(E[S_T]-F).
\]

$M(k)$는 단순히 $E[S_T]P(S_T\leq k)$가 아니다. 환율이 낮은 상태의 평균이 전체 평균보다 낮다는 점을 적분이 반영한다.\par

\par\medskip\noindent\textbf{심화: $M(a)$ 식의 유도}\quad $X=\ln(S_T/S_0)$의 밀도에 $S_0e^x$를 곱해 $x\leq\ln(a/S_0)$에서 적분한다. 지수부 $x-(x-\mu)^2/(2\sigma^2)$를 $-(x-\mu-\sigma^2)^2/(2\sigma^2)+\mu+\sigma^2/2$로 완전제곱하면 앞의 계수 $S_0e^{\mu+\sigma^2/2}$와 표준정규 누적분포 $\Phi$가 나온다.\par

\subsection{슬라이드 18 · 비선형 CFaR는 원래 환율의 5\% 분위수가 아닐 수 있다}
이 장에서 $CFaR_{5\%}$는 \textbf{현금흐름 분포의 하위 5\% 분위수} $q_{0.05}(CF_Q)$로 정의한다. 선도 후 각 구간의 현금흐름은 $h_j(s)=-150+(a_j+Q)s-QF$다. 전체 누적분포는\par

\[
P(CF_Q\leq c)=\sum_{j\in\{\ell,m,h\}}
P\bigl(S_T\in I_j,\ h_j(S_T)\leq c\bigr).
\]

$Q\geq-45{,}000$이고 낮은 환율 구간의 매우 희귀한 점프를 따로 세면 하위 5\%는 중간 구간의 $s_{0.05}\approx0.007164$에서 형성된다. 하지만 $Q<-45{,}000$이면 \textbf{중간 구간의 기울기 $45{,}000+Q$가 음수}가 되어, 중간 구간에서는 더 높은 환율이 오히려 더 낮은 현금흐름을 만든다. 기존의 $s_{0.05}$를 기계적으로 대입하면 잘못된 값이다.\par

\subsection{슬라이드 19 · 기울기가 뒤집힌 뒤의 분위수 찾기}
$-55{,}000<Q<-45{,}000$라면 높은 환율 구간의 기울기는 여전히 양수이고 중간 구간은 음수다. 매우 낮은 환율 구간의 확률 $p_\ell=P(S_T\leq k_2)\approx0.000165$를 아래쪽 꼬리로 포함한다. 중간 구간에서 현금흐름이 가장 낮은 쪽은 $S_T$가 $k_1$에 가까운 쪽이므로, 그 안에서 추가로 필요한 확률은 $0.05-p_\ell$이다. 따라서 임계 환율 $s^*$는\par

\[
P(s^*\leq S_T\leq k_1)=0.05-p_\ell,
\quad P(S_T\leq s^*)=P(S_T\leq k_1)-(0.05-p_\ell)
\approx0.183283.
\]

정규 누적분포를 역으로 사용하면 $s^*=S_0\exp\{\mu+\sigma\Phi^{-1}(0.183283)\}\approx0.007435$다. 아주 낮은 구간에서 $h_\ell(s)$가 다시 위로 올라가는 극단적 꼬리는 이 수치 예제에서 무시할 만큼 작은 확률이다. 정확한 일반 문제라면 위 누적분포 식으로 \textbf{모든 구간의 역상}을 확인해야 한다.\par

\subsection{슬라이드 20 · 비선형 목표 200을 맞추는 수량}
중간 구간에서 $CFaR_{5\%}=h_m(s^*)=-150+(45{,}000+Q)s^*-QF$다. $F=0.007720$ USD/JPY, $s^*\approx0.007434919$를 넣고 목표 $200$에 맞추면\par

\[
Q=\frac{200-[-150+45{,}000s^*]}{s^*-F}
\approx-54{,}120\quad\text{(백만 엔)}.
\]

이때 기대 현금흐름은 앞 평균식으로 약 $257.637$백만 달러다. 선도환율의 다섯째·여섯째 소수자리 차이가 $Q$를 크게 바꿀 수 있다. 예컨대 같은 $s^*$와 목표에서 $Q=-53{,}131$을 얻으려면 $F\approx0.0077253$을 사용해야 한다. 입력 환율과 헤지수량을 함께 표기하고 검산해야 하는 이유다.\par

\subsection{슬라이드 21 · 평균-하방 곡선에 꺾임이 생기는 이유}
선도 수량을 0에서 더 음수로 움직이면 $E[CF_Q]$는 $Q$의 일차식으로 변한다. 그러나 $CFaR$는 하위 5\%를 만드는 환율 상태가 바뀌는 지점에서 다른 식으로 이어진다. 그래서 평균-하방 그래프에 꺾임이 생긴다. $Q=-45{,}000$은 중간 구간의 기울기를 0으로 만드는 경계이고, 그 왼쪽에서는 목표 현금흐름을 높이기 위해 \textbf{환율의 다른 분위수}를 사용해야 한다. 그래프의 한 점만 보고 ‘완전헤지’라 부르면 이 구조를 놓친다.\par

\subsection{슬라이드 22 · 선형 선도의 한계와 위험지표의 한계}
선형 사업 현금흐름 $A S_T+B$에는 $Q=-A$가 모든 환율 상태에서 같은 현금흐름을 만든다. 비선형 사업의 $a(S_T)S_T+B$에는 구간마다 다른 기울기와 점프가 있어 한 개의 상수 $Q$로 모든 상태를 평평하게 만들 수 없다. 또한 VaR나 CFaR는 \textbf{한 분위수}만 알려 준다. 같은 5\% 분위수를 가진 두 헤지가 1\% 분위수, 평균, 최악 시나리오, 급전 필요액에서는 크게 다를 수 있다.\par

이를 보완하려면 여러 분위수와 스트레스 시나리오, 평균 초과손실 또는 하위 꼬리의 평균인 기대부족액(ES)을 함께 비교한다. 사업 수요의 임계값 자체가 불확실하면 그 모형도 민감도 분석해야 한다.\par

\subsection{슬라이드 23 · 비용 0의 계약에도 경제적 비용은 있다}
공정 선도는 처음의 계약가치가 0일 수 있지만 원하는 환율 방향에서 벌 수 있었던 이익을 포기한다. $E[S_T]>F$인 예에서 환율 상승에 유리한 기업이 엔화를 선도 매도하면 기대 현금흐름이 내려간다. 이것이 위험 감소와 맞바꾼 \textbf{암묵적 비용}이다. 반면 선물이라면 증거금 조달의 현금 비용도 생길 수 있다. 기대값이 줄어드는 정도는 계약의 ‘수수료’와 구별해야 한다.\par

\subsection{슬라이드 24 · 베이시스·교차헤지로 확장하기}
실제 기업이 받는 환율과 선물의 정산환율, 현물의 품질과 선물의 인도 규격, 사업의 현금흐름 날짜와 선물 만기가 다르면 완전한 상쇄가 되지 않는다. 일반적으로 $\Delta CF\approx A\Delta S+Q\Delta F$일 때 분산 최소화 수량은\par

\[
Q^*=-A\frac{\operatorname{Cov}(\Delta S,\Delta F)}{\operatorname{Var}(\Delta F)}.
\]

분산 $\operatorname{Var}(A\Delta S+Q\Delta F)$를 $Q$로 미분해 $2A\operatorname{Cov}(\Delta S,\Delta F)+2Q\operatorname{Var}(\Delta F)=0$으로 놓으면 얻는다. $\Delta F=\Delta S$일 때 $Q^*=-A$라는 완전헤지로 돌아간다. \textbf{상관관계가 1보다 낮으면 남는 베이시스 위험}을 다음 단계의 출발점으로 삼는다.\par

\end{document}
